\documentclass[11pt,a4paper]{article}
\usepackage[T1]{fontenc}
\usepackage[utf8]{inputenc}
\usepackage{lmodern}
\usepackage[margin=25mm,headheight=14pt]{geometry}
\usepackage{amsmath,amssymb,amsthm,mathtools}
\usepackage{booktabs,longtable,array}
\usepackage{microtype}
\usepackage{xcolor}
\usepackage{xurl}
\usepackage[unicode,colorlinks=true,linkcolor=blue!45!black,citecolor=blue!45!black,urlcolor=blue!45!black]{hyperref}
\usepackage{fancyhdr}
\pagestyle{fancy}
\fancyhf{}
\fancyhead[L]{\small CX\_GRH: exact hypotheses and transfers}
\fancyhead[R]{\small 21 November 2026}
\fancyfoot[C]{\thepage}
\setlength{\parskip}{4pt}
\setlength{\parindent}{0pt}
\setlength{\emergencystretch}{2em}
\setcounter{tocdepth}{2}
\title{\textbf{CX\_GRH}\\[5pt]
\large What the stipulated Dirichlet GRH and ETH theorems,\newline
and the conditional all-number-field Hecke claim,\newline
change in computational complexity}
\author{Research run CX\_GRH-20261121}
\date{21 November 2026 (Europe/Berlin)}
\begin{document}
\maketitle
\begin{abstract}
We identify the analytic premise at the point where it enters each
selected algorithm and then prove the corresponding transfer.
The stipulated Dirichlet theorem covers the Miller--Bach test,
prime-field nonresidues, and qualified primitive-root and
extension-field constructions. Koiran's general AM proof uses
Dedekind zeta functions of bounded, possibly nonnormal solution fields;
it is covered by the all-number-field Hecke claim, conditionally on
that claim. We preserve the equation-access and output-size distinctions
in GCT~V, the degree dependence of finite-field factoring, and the
separate search costs in class-group computation. The supplied ETH
theorem refutes the usual exponential-time hypotheses and gives
explicit parameter-preserving counting transfers and its stated
nearly linear Boolean circuit lower bound. These deductions supply
no P/NP or VP/VNP separation or collapse. A separate part carries
every finding into schemes over $\mathbb Z$, finite étale covers,
representation $K_0$, exterior and Schur operations, Hilbert schemes,
and the appropriate weight and Artin-motive language. Exact source
locators, inspected archive scans, proof checks, and unresolved
extensions are recorded explicitly.
\end{abstract}
\clearpage
\tableofcontents
\clearpage
\part{Checked analytic and computational findings}


\section{Conventions, premises, and the principal conclusions}
\label{sec:premises}

This paper answers the supplied \texttt{CX\_GRH} brief by auditing the
\emph{analytic hypothesis of each proof}, then substituting the two stipulated
2026 statements at precisely that point.  A proved transfer is different from
a new proof of its cited analytic or complexity theorem.  Our finding IDs
identify transfers, encoding refinements, exact dictionaries, and refutations
of specific proposed shortcuts.  They do not claim discovery of the classical
algorithms or of the lower bound already stated in the supplied ETH preprint.

All computational bounds use the deterministic Turing bit model unless a
source's arithmetic model is expressly retained.  A number field supplied by
an algebraic point need not be normal.  A character of a ray class group of a
number field is a Hecke character of that field; the phrase ``Dirichlet
character'' in an older paper does not automatically mean a character over
$\mathbb Q$.  We use $\log$ for the natural logarithm except where explicitly
subscripted.  A finite number of small inputs can always be handled directly.

\subsection{The stipulated theorems and the conditional draft claim}

\paragraph{The theorem $A_D$.}
As instructed, we cite Alp\"oge's Dirichlet GRH theorem as a theorem.
The exact supplied interfaces are
\texttt{DirichletLNonvanishingRightOfHalf} and
\texttt{DirichletLZerosOnLine} in \texttt{ALPOGE\_Statements.lean}, proved by
\nolinkurl{LFunction_ne_zero_of_half_lt_re} and
\nolinkurl{LFunction_eq_zero_re} in \texttt{ALPOGE\_Main.lean}.
The conductor-one case includes the Riemann zeta function; the latter file
also explicitly states \texttt{riemannHypothesis}.  The narrative conclusion
is \cite[Remark, p.~4, step~20]{Alpoge2026}.  The theta, positivity, period,
and cohomology steps~2--19 describe the proposed proof's method; none is an
extra input to our algorithmic transfers.  We infer no new standard-conjecture
or purity statement merely from that vocabulary.

\paragraph{The claim $M$.}
This denotes Masullo's draft claim \cite{Masullo2026} of GRH for Hecke $L$-functions over every
number field in the usual unitary normalization.  Only finite-order characters,
including the trivial character, are needed here.  The claim is stated in
Chapter~1, \S1.2, especially \S1.2.3 (PDF pp.~12--16);
the conventions are Chapter~3, Definition~3.2.1.2,
Proposition~3.2.1.3, and Definition~3.2.1.4 (PDF pp.~123--124).
Theorem~12.5.1.1 (pp.~519--520) and Theorem~13.3.1.1 (p.~522) are
trace-formula antecedents, not themselves a substitute for the GRH claim.
The displayed heading Theorem~13.3.1.2 on PDF p.~522 has no theorem body;
its footnote reads ``Insert statement.''  The preceding sentence has an
unresolved theorem reference.  We therefore cite the introductory claim $M$,
not an invented completed statement of Theorem~13.3.1.2.
This is a correction of the locator, not an authenticity or validity audit.
Every conclusion using $M$ is marked \textbf{conditional on $M$}.

\paragraph{The theorem $E$.}
The supplied OpenAI preprint's Theorem~1.1, source label
\texttt{thm:main}, is accepted as instructed.  For each fixed $k\geq2$, it
counts satisfying assignments of $k$-CNF formulas on $n\geq20$ variables and
$m$ clauses in polynomial space and time
\[
 (n+m)^{O_k(1)}
 2^{O_k(n\log\log n/\sqrt{\log n})}.
\]
For each fixed density $D$, when $m\leq Dn$ the exponential factor improves
to $2^{O_{k,D}(n/\sqrt{\log n})}$.  These are the explicit source bounds,
which imply the asymptotic form in the brief.  No uniformity in a growing
clause width $k(n)$ is assumed.  All consequences of $E$ in this paper are
independent of $A_D$ and $M$.  In the component proofs below,
$H_A$ and $\mathsf D$ are aliases for $A_D$, and $H_M$ and
$\mathsf H$ are aliases for $M$; they add no hypotheses.

\paragraph{What was supplied and checked.}
We read the two Lean interface files and the supplied mathematical sources.
The full Lean development and the claimed 175,187-file scan TSV are not in
the attachment.  We consequently make no claim of replaying that development,
checking its kernel, or inspecting that scan.  This does not change the
brief's instruction to accept the cited theorems.  Our new proofs have been
read in full by their author and checked by a second agent or the coordinator;
the event register distinguishes those checks from cited background and
from formal verification.  No new proof in this paper is advertised as
formally verified.

\subsection{Outcome map}

The complete versions, input promises, locators, and proofs occur below.
``Via $M$'' always means conditional on that draft claim, with no further
GRH premise.  ``Not discharged by $A_D$'' reports the identified proof gate,
not logical independence of analytic conjectures.

\begin{longtable}{@{}p{0.27\textwidth}p{0.37\textwidth}p{0.29\textwidth}@{}}
\toprule
Problem or conclusion & Exact gate of the cited route & Status and qualification\\
\midrule
\endhead
Miller test with Bach cutoff & Quadratic Dirichlet characters over $\mathbb Q$ & Via $A_D$; deterministic polynomial time for this test. PRIMES was already in P.\\
Prime-field power nonresidues; primitive roots & Rational Dirichlet characters & Via $A_D$; certified primitive-root construction retains supplied factorization of $p-1$, or report a candidate list.\\
Hilbert Nullstellensatz in AM & $\zeta_E$ for bounded algebraic solution fields $E$, not necessarily normal & Via $M$: $\mathrm{HN}\in\mathrm{AM}\subseteq\Pi_2^P$. A bounded abelian-witness promise is covered by $A_D$.\\
Equation-given homogeneous parameters & The same Koiran gate & Via $M$: the standard-graded search lies in $\mathrm{FP}^{\Sigma_2^P}$; unconditional functional PSPACE remains available.\\
Succinct explicit varieties, strict small separating systems & Koiran in the historical GCT~V EXPH proof & Via $M$ for that proof; construction already has an unconditional PSPACE theorem in later literature.\\
Succinct explicit varieties, optimal-cardinality h.s.o.p. & Koiran plus equations computable in singly exponential time & Via $M$ only with the additional equation-access premise retained.\\
Class-group factor bases and verification & Class/ray Hecke characters over $K$; $\zeta_{H_{K,\mathfrak m}}$ suffices & Via $M$; relation-search and smoothness hypotheses are separate.\\
General finite-field factoring & $\zeta_{\mathbb Q}$ and $\zeta_{\mathbb Q(\zeta_{2\nu},p^{1/\nu})}$ for the explicit construction degrees $\nu\le d^2$ & Via $M$: the complete proof CX-N15 gives $(d^{\log d}\log q)^{O(1)}$ time.\\
Extension-field nonresidues & Rational Dirichlet characters when the specified prime $r$ divides $p-1$ & Via $A_D$ in every supplied extension model; CX-N13. Other selected algorithms retain their stated character families.\\
ETH and SETH families & Only $E$, with fixed-width quantifiers & Their usual deterministic lower-bound hypotheses are false; a false premise gives no blanket target-problem algorithm.\\
Circuit consequence & $E$ plus the cited algorithm-to-lower-bound theorem & Some $L\in\mathrm E^{\mathrm{NP}}$ has $C_L(n)>c n\sqrt{\log n}$ infinitely often.\\
VP, VNP, permanent, GCT & Symbolic arithmetic circuits and equivariant coordinate data & Numerical counting transfers are proved; no separation or collapse of VP/VNP or P/NP is obtained by these deductions.\\
\bottomrule
\end{longtable}

\subsection{CX-R01: the exact inherited number-field scope}
\label{finding:R01}

\textbf{Finding.} The theorem $A_D$ implies RH for $\zeta_K$ for every finite
abelian extension $K/\mathbb Q$.  Conditional on $M$, RH holds for $\zeta_K$
for every number field and for every finite-order Hecke character of every
number field.  No Artin-holomorphy assertion is needed for these transfers.

\begin{proof}
For an abelian $K/\mathbb Q$, use the primitive-character factorization
\[
 \zeta_K(s)=\prod_{\chi\in\widehat{\operatorname{Gal}(K/\mathbb Q)}}L(s,\chi)
\]
with each character replaced by its associated primitive Dirichlet character;
see \cite[Theorem~4.3, p.~33]{Washington1997}.  This conductor convention
includes the correct ramified Euler factors.  Every nontrivial factor is
entire and the trivial factor has its sole pole at $s=1$.  In the open strip
$0<\Re s<1$ off the central line all factors are finite and nonzero by $A_D$.
Their finite product is nonzero there, proving the assertion.

For arbitrary $K$, the finite Hecke $L$-function of the trivial idele class
character is $\zeta_K$.  Archimedean gamma factors and the conductor
exponential factor are holomorphic and nonzero in that open strip.  Thus
their inclusion in a completed $L$-function does not change the strip zeros.
The claim $M$ places those zeros on the central line.  Every finite-order
ray-class character is unitary and is included in $M$.  If an algorithm uses
an imprimitive Euler product, the removed local factors have form
$1-\alpha N\mathfrak p^{-s}$ with $|\alpha|=1$; their zeros have real part
zero, so they do not change this open-strip conclusion.
\end{proof}

\subsection{CX-R02: a concrete obstruction to the rank-one shortcut}
\label{finding:R02}

\textbf{Finding.} For $K=\mathbb Q(\sqrt[3]{2})$, the permutation
representation on $\operatorname{Hom}_{\mathbb Q}(K,\overline{\mathbb Q})$
is not a virtual integral sum of rank-one $G_{\mathbb Q}$-representations.
In particular the proposed representation-theoretic factorization of every
solution-field zeta function into rational Dirichlet factors fails.

\begin{proof}
The polynomial $x^3-2$ is Eisenstein at $2$, and its discriminant $-108$ is
not a rational square.  Its normal closure therefore has group $S_3$.
The permutation representation $\rho$ has dimension $3$, trace $3$ at the
identity, and trace $0$ at a three-cycle.  Choose $u,v\in G_{\mathbb Q}$
mapping to $(12)$ and $(23)$.  Their commutator $c$ maps to a three-cycle.
Every rank-one character $\chi$ of $G_{\mathbb Q}$, whether or not it
factors through this $S_3$ quotient, satisfies $\chi(c)=\chi(1)=1$.
Consequently every integral virtual sum of rank-one characters has trace
at $c$ equal to its dimension.  The representation $\rho$ has trace $0$
there and dimension $3$.  Trace is additive in exact sequences, so equality
in the representation Grothendieck group is impossible.
\end{proof}

This refutes a specified dictionary, not a logical independence statement
$A_D\nRightarrow\mathrm{RH}(\zeta_K)$ and not every possible future analytic
reduction.  The related failure of replacing an arbitrary solution by an
abelian-valued solution is proved separately in CX-K04.  Part~II states
these findings again in the language of finite étale covers and $K_0$.



% CX_GRH: standalone section fragment. No preamble.
% Bibliography keys are supplied in koiran_bibliography.tex.
% Checked 2026-11-20 UTC by /root/koiran_gct, GPT-6 Astra Pro,
% inherited effort (setting not exposed). Independent full-proof checks:
% /root/eth_consequences and /root, with qualifications incorporated.

\section{Koiran and Noether normalization: the exact analytic gate}

The analytic hypothesis needed by Koiran's proof is RH for the Dedekind
zeta functions of bounded solution fields. These fields can be arbitrary
number fields. They need not be Galois over $\mathbb Q$, and passing to a
Galois closure is unnecessary. Masullo's all-number-field Hecke-GRH claim,
if accepted as $H_M$, supplies this hypothesis by its trivial-character
case. The stipulated Alp\"oge Dirichlet-GRH theorem $H_A$ supplies the
corresponding input for abelian extensions of $\mathbb Q$, but the checked
factorization argument does not extend this to all solution fields.

Write
\[
 H_D:\quad \text{RH holds for $\zeta_E$ for every number field $E$.}
\]
Throughout this section, $H_M$ and $H_A$ have the precise meanings fixed
in the premise register. In particular, $H_M$ remains a conditional
claim. Our conclusions under $H_M$ are transfers of established
complexity theorems, together with the encoding and scope refinements
proved below. They are not independent unconditional proofs of those
theorems.

\subsection{Source conventions, numerical locators, and representations}

Koiran's published paper and extended report require separate locators.
The journal version \cite[\S2, Theorem~2, pp.~276--277]{K96} displays
$\mathrm{RP}^{\mathrm{NP}}$, followed by containment in
$\Pi_2^{\mathrm P}$. The extended DIMACS report explicitly displays
$\mathrm{AM}\subseteq\Pi_2^{\mathrm P}$ in \S2, Theorem~2,
printed pp.~5--6 \cite{K96TR}. Its protocol uses public randomness before
the nondeterministic collision witnesses. We use this report for the
requested AM formulation. Koiran's 8 August 1997 erratum, first page,
second bullet, corrects the lower bound in Theorem~8 and states that
Theorem~1 and the remainder of the paper are unaffected
\cite{K97Err}. All transfers here use the corrected version.

The input in Koiran's \S1.1 is a finite list of integer polynomials in
the Turing model, with sparse encoding. Theorems~4 and~7, Lemma~3, and
Theorem~8 supply the bounded algebraic witness and reduction modulo
primes. Theorem~9 and Corollary~1 in \S5, journal p.~282 and DIMACS
printed p.~12, isolate the effective prime-ideal theorem. The relevant
analytic result is Lagarias--Odlyzko \cite[Theorem~1.1]{LO77}, whose
hypothesis is RH for the Dedekind zeta function of the top field.

For GCT~V \cite{GCT5}, the attached \texttt{main.tex} labels and the
numerical statements of arXiv:1209.5993v5 agree as follows. Page numbers
in this table refer to that arXiv PDF, not the journal pagination.

\begin{center}
\small
\begin{tabular}{@{}p{.14\linewidth}p{.25\linewidth}p{.43\linewidth}r@{}}
\hline
Statement & Attached source label & Use & Page \\
\hline
Theorem 2.8 & \texttt{tstrongexp} & Strengthened black-box PIT, EXPH under GRH & 24--25 \\
Proposition 2.9 & \texttt{pstandard} & Standard black-box PIT, PH under GRH & 24 \\
Theorem 2.10 & \texttt{tkoiran} & Circuit and homogeneous HN & 24 \\
Lemma 3.1 & \texttt{lnnl0} & Geometric normalization & 26 \\
Lemma 3.2 & \texttt{lnoether20} & Graded integrality criterion & 26 \\
Theorem 3.4 & \texttt{tnoetherph} & Equation-given h.s.o.p., functional PH & 27--28 \\
Theorem 4.1 & \texttt{tnnldetexpspace} & Determinant h.s.o.p., extra equation hypothesis & 30 \\
Theorem 4.10 & \texttt{tssopverifydet} & Determinant s.s.o.p., construction/verification & 34 \\
Definition 5.1 & \texttt{dexpvariety} & Explicit families and succinct circuits & 35--36 \\
Theorem 5.5 & \texttt{tnnlexplicithsop} & Explicit h.s.o.p., extra equation hypothesis & 39 \\
Theorem 5.11(a) & \texttt{tnnlexpfull} & Border hitting sets give strict separating systems & 43 \\
Theorem 5.12 & \texttt{tssopexplicit} & Explicit strict s.s.o.p., construction/verification & 44 \\
Theorem 7.3 & \texttt{thsop} & Simultaneous-conjugation invariant ring & 50 \\
\hline
\end{tabular}
\end{center}

An h.s.o.p. has exactly $\dim R$ algebraically independent homogeneous
parameters and makes $R$ finite over their polynomial subring. A small
system of parameters may have polynomially many elements, with no
optimal-cardinality requirement. The succinct small systems in GCT~V
are specified through small circuits and are a different output regime
from explicitly listing exponentially many ambient coordinates.
``Functional PH'' means a construction with polynomial-time oracle
access to a fixed level of PH. It is not a decision-language assertion
and does not mean ordinary polynomial-time construction.

The standard specification preceding GCT~V Theorem~3.4 is a homogeneous
affine cone defined by homogeneous polynomial equations represented by
arithmetic circuits with rational constants. The characteristic-zero
statements here use this effective Turing representation. We do not
assert Turing bounds for arbitrarily presented transcendental
coefficients. The stronger refinement in CX-K05 explicitly uses
standard grading. GCT~V Definition~5.1 also allows differing positive
weights; the cited mixed-weight theorem is kept separate in CX-K06.

More precisely, Definition~5.1 presents
$W_n=\overline{\psi_n(K^{r_n})}\subseteq K^{m_n}$ with
$r_n=\operatorname{poly}(n)$, $m_n=n^{\Omega(1)}$, and
$\log m_n=\operatorname{poly}(n)$. Thus
$\dim W_n\le r_n$, although its ambient dimension may be exponential.
Writing $\psi_n=(f_1,\ldots,f_{m_n})$, the succinct input is a
uniformly constructible rational circuit, of polynomial total bit
length, for $F_n(v,x)=\sum_j f_j(v)g_j(x)$. The $g_j$ are linearly
independent; the $f_j$ and $g_j$ are homogeneous of polynomial degree,
and $F_n$ has polynomial total degree. Definition~5.6 measures a short output by circuits for
$\psi_n^*(s)$, including their rational constants' bit lengths.
Definition~5.7 calls an element strict when it is a specialization
of a fixed homogeneous-in-$v$ component of $F_n$ at a polynomial-bit
integer point $x=b$. Hence a strict output is encoded by polynomially
many such points and component indices, without listing $m_n$
ambient coefficients. The dimension bound alone does not require the
small system's cardinality to equal the dimension.

\subsection{CX-K01: the root-field gate without a Galois closure}
\label{cx:k01}

\paragraph{Statement.}
In Koiran's proof, the analytic premise can be localized inputwise to
RH for $\zeta_E$, where $E$ is the field of a bounded algebraic solution
chosen using Koiran's Theorems~4 and~7. The uniform family of possible
$E$ is unrestricted. Neither Artin holomorphy nor passage to a splitting
field is required. If one uses $\operatorname{Li}(x)$ rather than
Koiran's comparison with the rational prime-counting function, no
separate RH-for-$\mathbb Q$ premise is needed at this step.

\begin{proof}
Let $S$ be a satisfiable integer polynomial system of encoding length
$N$. Koiran's Theorems~7 and~4 give a solution
$a=(a_1,\ldots,a_t)$, a primitive element $\theta$ for
$E=\mathbb Q(a_1,\ldots,a_t)$, an integer minimal polynomial $R$ of
$\theta$, and integer denominators in expressions for the $a_i$.
The degree, logarithmic discriminant bounds, and denominator logarithms
are all at most $2^{q(N)}$ for a fixed polynomial $q$.

We can make the primitive element integral without losing that scale.
If $R$ has leading coefficient $c$ and degree $d$, then
$c^{d-1}R(T/c)$ is a monic integral polynomial for
$\beta=c\theta$. Thus $\mathbb Z[\beta]\subseteq\mathcal O_E$
has finite index, the field discriminant satisfies the usual
polynomial-discriminant/index relation, and all $a_i$ lie in
$\mathcal O_E[1/B]$ for an integer $B\ne0$ with singly exponentially
bounded logarithm. Put $J=[\mathcal O_E:\mathbb Z[\beta]]$ and
$D_E=\operatorname{disc}(E)$. The logarithm of $J$ has the same bound.

Apply Lagarias--Odlyzko Theorem~1.1 to the trivial extension $E/E$.
It is a Galois extension, regardless of whether $E/\mathbb Q$ is
normal. Under RH for $\zeta_E$, the number $\pi_E(x)$ of prime ideals
of $\mathcal O_E$ of norm at most $x$ satisfies
\begin{equation}
 \pi_E(x)=\operatorname{Li}(x)+
 O\!\left(\sqrt{x}\log(|D_E|x^d)+\log |D_E|\right),
 \qquad d=[E:\mathbb Q],
 \label{eq:k01primeideal}
\end{equation}
with an effective absolute implied constant. This is an invocation of
the cited analytic theorem, not a new proof of that theorem.

A prime ideal of residue degree at least two and norm at most $x$
lies above a rational prime at most $\sqrt{x}$. At most $d$ prime
ideals lie above any one rational prime. Hence these higher-degree
prime ideals contribute at most $d\pi_{\mathbb Q}(\sqrt{x})\le
d\sqrt{x}$. Discard also the prime ideals above divisors of $BJ$;
there are at most $d\omega(BJ)$, where $\omega$ counts distinct prime
divisors. At any remaining residue-degree-one prime ideal
$\mathfrak p$, reduction maps the solution coordinates into
$\mathcal O_E/\mathfrak p=\mathbb F_p$. The equations of $S$ remain
valid. Since each rational prime has at most $d$ prime ideals above it,
the number of rational primes $p\le x$ for which $S$ is soluble is at
least
\begin{equation}
 \frac{1}{d}\left[
 \operatorname{Li}(x)
 -O\!\left(\sqrt{x}\log(|D_E|x^d)+d\sqrt{x}+\log|D_E|\right)
 -d\omega(BJ)\right].
 \label{eq:k01usefulprimes}
\end{equation}

All quantities outside $x$ and $\log x$ in this expression are bounded
by $2^{\operatorname{poly}(N)}$. Choose
$x=2^{P(N)}$ for a sufficiently large fixed polynomial $P$.
The main term in \eqref{eq:k01usefulprimes} then dominates its errors
and the unconditional inconsistent-system bound in Koiran's
Theorem~5 by the gap required in his Theorem~1. Its primes still have
polynomial bit length. Koiran's unchanged public-coin protocol in
DIMACS Theorem~2 therefore applies. This identifies the exact analytic
gate and shows that no Galois closure is needed. In particular, the
potential factorial degree growth of a splitting field is avoided.

Finally, every number field $E$ has a primitive element $\theta$.
Clear denominators in its irreducible polynomial $R$. All roots of
$R$ generate conjugate copies of $E$. Thus the single equation
$R(x)=0$ realizes this isomorphism type among the possible solution
fields. The input family is therefore not confined to abelian fields.
This observation concerns the fields occurring in the checked proof;
it does not establish that the complexity conclusion logically
requires $H_D$.
\end{proof}

The weaker unconditional prime-ideal estimate in
Lagarias--Odlyzko Theorem~1.3 starts in the range
$x\ge\exp(10d(\log|D_E|)^2)$. With singly exponential $d$ and
$\log|D_E|$, that range need not preserve polynomial bit length of
the primes. It therefore does not substitute into this proof with the
same complexity bound. This is a limitation of the checked route,
not an independence claim about Dirichlet GRH and $\mathrm{HN}\in
\mathrm{AM}$.

\subsection{CX-K02: conditional HN for circuits and nonzero homogeneous points}
\label{cx:k02}

\paragraph{Statement.}
Under $H_M$, feasibility over $\mathbb C$ of polynomial systems given
by arithmetic circuits with rational constants belongs to
\[
 \mathrm{AM}\subseteq\Pi_2^{\mathrm P}.
\]
The same bound holds for existence of a nonzero solution of a
homogeneous such system. The complementary infeasibility predicates
belong to $\Sigma_2^{\mathrm P}$.

\begin{proof}
For every number field $E$, the Hecke $L$-function of the trivial
character is $\zeta_E$. Consequently $H_M\Rightarrow H_D$.
Apply the corrected Koiran DIMACS Theorem~2, or CX-K01 followed by its
unchanged hashing protocol. This gives the asserted AM bound for
Koiran's sparse input representation.

For a circuit representation, introduce a variable $y_g$ for each
gate. A sum gate and a product gate give, respectively,
\[
 y_g-y_h-y_j=0,\qquad y_g-y_hy_j=0.
\]
A rational constant $a/b$ gives $by_g-a=0$, with $b\ne0$.
Wire the input gates to the original variables, and set each output
gate variable to zero. The result is a polynomial-size sparse system
of degree at most two and polynomial total coefficient bit length.
Every original solution has exactly one extension to gate values.
Conversely, induction over the circuit forces all gate variables to
equal the intended evaluations. This is a deterministic
polynomial-time many-one reduction, which preserves AM membership.

For a homogeneous system in $x_1,\ldots,x_t$, adjoin variables
$u_1,\ldots,u_t$ and the equation
\[
 \sum_{i=1}^t u_i x_i=1.
\]
A nonzero solution admits such a $u$ by inverting a nonzero coordinate,
and a solution to this enlarged system necessarily has $x\ne0$.
This again is a polynomial-size circuit reduction. The standard
containment used in Koiran is $\mathrm{AM}\subseteq\Pi_2^{\mathrm P}$;
complementing gives $\Sigma_2^{\mathrm P}$. No containment
$\mathrm{AM}\subseteq\Sigma_2^{\mathrm P}$ is asserted.
\end{proof}

\subsection{CX-K03: the bounded abelian-witness consequence of the stipulated Dirichlet theorem}
\label{cx:k03}

\paragraph{Statement.}
Fix an explicit polynomial $P$. Consider the promise problem whose
no-instances are all inconsistent rational circuit systems and whose
yes-instances have a solution as in CX-K01 with a finite abelian
Galois field $E/\mathbb Q$ and degree, logarithmic defining-polynomial
discriminant, and denominator logarithms at most $2^{P(N)}$.
Under $H_A$, this promise problem is in AM. Its verifier need not
compute $E$, certify abelianness, or recognize the promise.

\begin{proof}
For finite abelian $E/\mathbb Q$, the Dedekind zeta function factors
as the product of primitive Dirichlet $L$-functions attached to
characters of $\operatorname{Gal}(E/\mathbb Q)$, with the conductor
convention giving the correct ramified Euler factors. The trivial
factor is $\zeta_{\mathbb Q}$. The stipulated $H_A$ therefore puts
all nontrivial zeros of $\zeta_E$ on the critical line.

On every promised yes-instance, CX-K01 now gives the useful-prime
bound with a uniform polynomial exponent. On every no-instance,
Koiran's Theorem~5 supplies the unconditional bound. The same gap
and public-coin protocol used in Koiran's DIMACS Theorem~2 apply.
The arithmetic promise proves completeness of the analytic estimate;
the protocol never checks it. Soundness on all inconsistent systems
is unchanged.
\end{proof}

The bound on the abelian witness is essential to the statement proved
here. Koiran's existence of a bounded solution does not say that a
solution satisfying an additional arithmetic property can be chosen
with that same bound. We therefore do not replace the quantitative
promise by the bare assertion that some abelian-valued solution exists.

\subsection{CX-K04: a counterexample to replacing all solution fields by abelian fields}
\label{cx:k04}

\paragraph{Statement.}
The satisfiable equation $x^3-2=0$ has no solution in any abelian
extension of $\mathbb Q$. Hence forcing every witness into an abelian
field is not an available general reduction.

\begin{proof}
Eisenstein's criterion at $2$ makes $x^3-2$ irreducible. Its real root
generates a degree-three subfield of $\mathbb R$ that omits the two
nonreal conjugates, so that root field is not normal over $\mathbb Q$.
The other two root fields are conjugate to it and are likewise
nonnormal. Every intermediate field of a finite abelian Galois
extension is Galois over $\mathbb Q$, since every subgroup of an
abelian group is normal. None of these root fields can be an
intermediate field. A root in an infinite abelian extension would
already lie in a finite abelian subextension, so that possibility
also fails.
\end{proof}

Only the abelian-witness shortcut is refuted. This counterexample does
not refute the possibility of a different proof of
$\mathrm{HN}\in\mathrm{AM}$ from $H_A$ alone.

\subsection{CX-K05: the functional PH level for equation-given h.s.o.p.'s}
\label{cx:k05}

\paragraph{Statement.}
Assume $H_M$. Let
$R=\mathbb Q[z_1,\ldots,z_t]/I$ be a nonzero standard positively
graded algebra, with $I$ specified by homogeneous rational-constant
arithmetic-circuit equations. Let $N$ be the total encoding length,
including an explicit list of the ambient coordinates, so $t\le N$.
An h.s.o.p. of rational linear forms and $\dim R$ can be constructed
in
\[
 \mathrm{FP}^{\Sigma_2^{\mathrm P}}
\]
with polynomial output bit length. These forms remain an h.s.o.p.
after extension of scalars to any characteristic-zero field.
Unconditionally the construction has a polynomial-space bound.

\begin{proof}
The bounded-coefficient part of GCT~V Theorem~3.4's proof gives a
computable polynomial $b(N)$ such that, for $d=\dim R$, there are
$d$ integral linear forms, each coefficient having at most $b(N)$
bits, whose common zero on the cone is only its vertex. The relevant
size argument in that proof bounds the equation degrees by $2^N$,
passes to a common homogeneous degree at most $2^{N^2}$, and uses a
nonzero resultant of degree with polynomial logarithm. In particular,
only a polynomial bit bound is needed for the coefficient grid;
enumerating all monomials is unnecessary. Nilpotents do not change
the set of geometric zeros or the dimension, so the same existence
bound applies to the algebra stated here.

For $L=(L_1,\ldots,L_s)$ let $\operatorname{Good}(I,L)$ mean that
the equations generating $I$, together with $L_1=\cdots=L_s=0$,
have no nonzero geometric solution. By CX-K02 this is a
$\Sigma_2^{\mathrm P}$ predicate. Fix a canonical binary encoding
of the bounded coefficient lists, padded to polynomial length. For
any prescribed prefix, the predicate
\[
 \exists L\;[\text{$L$ extends the prefix and }
                    \operatorname{Good}(I,L)]
\]
also belongs to $\Sigma_2^{\mathrm P}$: the bounded existential
guess of $L$ merges with the leading existential block. In explicit
quantifier form, if Good is $\exists u\,\forall v\,T(I,L,u,v)$
for polynomial-time $T$, the prefix predicate is
$\exists (L,u)\,\forall v\,[\operatorname{Prefix}(L)\land
T(I,L,u,v)]$.

Ask this predicate for $s=0,1,\ldots,t$. The least successful $s$
equals $d$. Existence at $d$ follows from the bounded-coefficient
argument. For $s<d$, the generalized principal ideal theorem gives
\[
 \dim R/(L_1,\ldots,L_s)\ge d-s>0,
\]
so a homogeneous cone cannot be supported only at the origin.
For $d=0$, the empty list works.

Once $s=d$ is known, recover a successful list one bit at a time.
Extend the current prefix by zero and ask if a successful completion
exists. Keep zero when the answer is yes, and otherwise keep one.
The invariant that the prefix has a completion holds initially and
is preserved by every step. There are polynomially many coefficient
bits and polynomially many dimension queries. This proves the
$\mathrm{FP}^{\Sigma_2^{\mathrm P}}$ bound, rather than merely a
language membership in PH.

We verify that the searched geometric condition is precisely the
required ring condition. Good makes $R/(L)$ finite-dimensional over
$\mathbb Q$: after extending to an algebraic closure, its spectrum is
supported at the origin, so its homogeneous irrelevant ideal is
nilpotent. Choose homogeneous representatives of a finite vector-space
basis of $R/(L)$. They generate $R$ over
$A=\mathbb Q[L_1,\ldots,L_d]$. Indeed, for any homogeneous element,
subtract a representative of its class modulo $(L)$. The difference
is a sum of terms $L_i r_i$, with each $r_i$ of strictly lower degree;
induction on degree proves the assertion. Thus $R$ is a finite
$A$-module and is integral over $A$.

Conversely, if $R$ is finite over $A$, its quotient by the images of
the $L_i$ is finite-dimensional over $\mathbb Q$, so its homogeneous
cone has no nonzero geometric point. Finally, integral extensions
preserve dimension, so $\dim A=d$. A nonzero kernel in
$\mathbb Q[T_1,\ldots,T_d]\to A$ would make $\dim A<d$.
Thus the parameters are algebraically independent and constitute an
h.s.o.p. This reasoning permits nilpotents and several irreducible
components.

All constructed coefficients are rational. Tensoring the finite-module
presentation and the injection of the parameter polynomial algebra
with a field extension preserves finiteness and injectivity. Hence
the same forms work after extending scalars.

For the unconditional bound, use the homogeneous-feasibility PSPACE
procedure cited in GCT~V Theorem~2.10. Enumerate bounded lists with
polynomial-space counters and test them using that procedure,
increasing $s$ until a successful list appears. All counters and
oracle workspaces are polynomial in $N$, giving a functional PSPACE
construction, consistent with GCT~V Theorem~3.4.
\end{proof}

The zero algebra is an excluded trivial case. The coordinate list
condition rules out encoding exponentially many unused free variables
by a binary dimension alone. Standard grading is needed for the
linear homogeneous output asserted here. The oracle bound does not
give an ordinary polynomial-time algorithm.

\subsection{CX-K06: the exact EXPH transfers for succinct explicit varieties}
\label{cx:k06}

\paragraph{Statement.}
Under $H_M$, GCT~V Theorem~5.12 gives construction and verification of
a strict succinct small system of parameters for an explicit family
$W_n$ as in Definition~5.1 in the functional exponential hierarchy;
the determinant version is Theorem~4.10. For an optimal-cardinality
h.s.o.p., the EXPH assertion in Theorem~5.5 retains the additional
hypothesis that defining equations can be computed in
$2^{\operatorname{poly}(n)}$ time; its determinant counterpart is
Theorem~4.1.

\begin{proof}
As in CX-K02, $H_M$ supplies $H_D$. Every application of Koiran in
GCT~V Theorem~2.8 concerns a finite integer or rational circuit system,
possibly of exponential encoding length. Its arithmetic premise is
therefore discharged by this same $H_D$.

For clarity about equation access, the proof of Theorem~2.8 itself
constructs an existential polynomial system of exponential total
encoding length and exponentially many variables for a candidate
border-circuit hitting set. Infeasibility of that system is the
hitting-set predicate. The construction uses
B\"urgisser \cite[Theorem~5.7]{Burgisser04}, as cited there; it does
not assume an exponentially small defining ideal for $W_n$.
The systems are constructed in deterministic exponential time.
Their good-candidate predicates have polynomial time and quantifier
bounds in the expanded length, with a constant number of
alternations. In terms of $n$, those bounds become exponential,
still with constantly many alternations. Searching a polynomial-length
description of a hitting set, followed by the transformation in
GCT~V Theorem~5.11(a), therefore stays in the functional exponential
hierarchy. The same predicates give the verification assertion in
Theorem~5.12. This proves the strict s.s.o.p. transfer and preserves
the output conventions of the cited theorem.

For the optimal-cardinality h.s.o.p. assertion, use GCT~V
Theorem~5.5 exactly as stated. Its proof invokes the argument of
Theorem~4.1 and hence Theorem~3.4. Their arithmetic premise is again
$H_D$. The extra equation algorithm supplies an expanded input of
length $L\le 2^{\operatorname{poly}(n)}$. Polynomial time and
polynomial quantifier length in $L$ become exponential time and
exponential quantifier length in $n$, while the number of alternations
stays constant. In the standard-graded case this substitution is
directly the algorithm in CX-K05. For differing positive weights,
we invoke the stated graded theorem with its representation
conventions; arbitrary linear combinations of generators of different
weights are not homogeneous.

Without equation access of the required time bound, this proof has
no first step of the stated size. The unconditional elimination
method discussed in GCT~V may produce double-exponential-length
equations. Supplying $H_D$ does not change that representation cost.
The extra assumption in Theorem~5.5 therefore remains part of this
specific EXPH conclusion.
\end{proof}

We use ``exponential hierarchy'' for the union of constant levels
defined by alternating exponential time with exponentially bounded
quantifiers; a functional version permits the corresponding oracle
construction. The argument proves membership at a fixed finite level
under the cited source conventions. It does not identify general
EXPSPACE with EXPH or claim that all normalizations require the
equation-access hypothesis.

\subsection{CX-K07: the existing unconditional PSPACE construction}
\label{cx:k07}

\paragraph{Statement.}
Polynomial-size succinct separating normalization systems for explicit
varieties already have unconditional polynomial-space constructions,
independent of $H_A$ and $H_M$. Thus their EXPH construction does not
become unconditional for the first time through either stipulated
2026 premise. Optimal-cardinality h.s.o.p.'s and verification of
arbitrary candidate systems remain different statements.

\begin{proof}
Forbes--Shpilka \cite[Theorem~7.1, preprint p.~21]{FS18} gives a
polynomial-space construction of a polynomial-size rational hitting
set, of polynomial bit complexity, for the closure of small
bounded-degree algebraic circuits over $\mathbb R$ or $\mathbb C$.
Its parameters are supplied in unary. Apply it to the polynomially
bounded parameters of the explicit family and feed the resulting set
to GCT~V Theorem~5.11(a). Polynomial-space construction followed by
the polynomial-time output transformation remains polynomial-space.

Guo--Saxena--Sinhababu \cite[Theorems~1.5 and~1.6, p.~6]{GSS19}
explicitly record the determinant and general explicit-variety
normalization consequences, and provide the corresponding finite-field
treatment. These results assume no GRH and predate the new premises.
Their outputs allow polynomial slack in the number of parameters;
they do not assert that the cardinality equals the dimension.
\end{proof}

This is a direct composition of cited algorithms and a correction to
the scope of the historical comparison. It is not a new discovery of
the PSPACE algorithm. Guo--Saxena--Sinhababu uses the term e.s.o.p.
while placing its construction in PSPACE. To avoid conflict with
GCT~V's polynomial-time convention for that term, the conclusion here
is stated as a polynomial-space construction of a polynomial-size
succinct separating system. Neither that theorem nor this composition
is an assertion about every arbitrary candidate-verification input.

\subsection{Remaining gates and verification status}

The full equation-given h.s.o.p. conclusion in
$\mathrm{FP}^{\Sigma_2^{\mathrm P}}$ remains conditional on $H_M$
in this run. From $H_A$ alone, CX-K03 establishes the explicitly
bounded abelian-witness promise version, and CX-K04 refutes the
unrestricted abelian-witness shortcut. No logical nonimplication
between $H_A$ and general HN membership is claimed. The succinct
optimal-cardinality h.s.o.p. EXPH transfer retains its equation-access
gate; the small separating-system construction already has the
unconditional PSPACE bound.

All proofs CX-K01--CX-K07 were read end to end after drafting by
\texttt{/root/koiran\_gct}, GPT-6 Astra Pro, on 20 November 2026 UTC;
the inherited effort setting was not exposed. The independent agents
\texttt{/root/eth\_consequences} and \texttt{/root} each read the
complete proofs. Their checks found no remaining substantive gap after
the standard-grading and output-scope qualifications were incorporated.
Koiran's protocol and erratum, the relevant GCT statements and displayed
proofs, and the later theorem statements were inspected. The full
historical analytic and border-hitting-set proofs are carried as cited
background, not claimed to have been independently reproved. Detailed
check events and source locators accompany the ledger.


\section{Number-theoretic algorithms: exact analytic gates}
\label{sec:number-algorithms}

Write \(\mathsf D\) for RH for all rational Dirichlet \(L\)-functions,
including \(\zeta_{\mathbb Q}\), and \(\mathsf H\) for the asserted RH for
Hecke \(L\)-functions over every number field. All character applications audited here are finite-order ones. The stipulated
Alp\"oge theorem discharges \(\mathsf D\). Every consequence of
\(\mathsf H\) below remains \emph{conditional on the all-number-field
claim of the Masullo draft}. A statement that a published proof's gate is
not discharged by \(\mathsf D\) is a scope audit, not a theorem of logical
nonimplication and not an exclusion of a different algorithm.

Finite fields are explicitly represented, for example as
\(\mathbb F_{p^m}=\mathbb F_p[T]/(g)\), with \(g\) supplied and irreducible.
A multiplicative primitive root generates \(\mathbb F_{p^m}^{\times}\);
a primitive element of the extension need only generate the field.
Integer factorizations supplied as input are stated explicitly.

\subsection{CX-N01: Miller testing with Bach's cutoff}
\label{cx-n01}

Handle bounded small integers and even inputs directly. For an odd
integer \(N\), write \(N-1=2^s u\), \(u\) odd, and test all
\[
 2\leq a\leq\min\{N-1,\lceil2(\log N)^2\rceil\}.
\]
Reject on a nontrivial gcd, or unless \(a^u\equiv1\pmod N\) or
\(a^{2^ju}\equiv-1\pmod N\) for some \(0\leq j<s\).
Under \(\mathsf D\), this is a deterministic polynomial-time primality test.

\paragraph{Exact source gate.}
Bach explicitly restricts \S6, pp.~369--373, to ordinary Dirichlet
characters; Theorem~2, p.~372, gives the subgroup bound with constant \(2\),
and p.~373 states the Miller/Solovay--Strassen witness consequence
\cite{CXNBach90}. Schoof's complete Theorem~2.2 proof, pp.~105--106,
explicitly states the smaller sufficient hypothesis of RH for
\emph{quadratic} Dirichlet characters \cite{CXNSchoof08}.

\paragraph{Full transfer proof.}
Those quadratic characters belong to \(\mathsf D\), so the cited criterion
gives a rejecting small base for every composite input. Every prime passes
all its nonzero bases. Increasing the strict source cutoff to the ceiling
cannot cause false acceptance, and bounded small inputs handle endpoint
conventions. There are \(O(b^2)\) bases for \(b=\lceil\log_2N\rceil\), each
requiring \(O(b)\) modular multiplications. The bit cost is
\(O(b^3M(b))\), where \(M(b)\) is multiplication cost; schoolbook arithmetic
gives \(O(b^5)\). Only the power of two in \(N-1\) is extracted.
This establishes this particular implementation; \(\mathrm{PRIMES}\)
was already in deterministic \(\mathrm P\) by AKS \cite[Theorem~4.1 p.~785 and Theorem~5.1 p.~789]{CXNAKS04}. \(\square\)

The entire Schoof proof was read for the hypothesis audit. Its squarefree
case uses quadratic characters modulo prime divisors and their products;
the repeated-prime case uses a smooth-number estimate and finite-group
counting. Miller's original locators are Theorem~2, p.~301, and the
Appendix, pp.~315--316 \cite{CXNMiller76}.

\subsection{CX-N02: a specified prime-field power nonresidue}
\label{cx-n02}

Given prime \(p\) and a specified prime \(r\mid p-1\), \(\mathsf D\) yields
an \(r\)-th power nonresidue in deterministic time polynomial in
\(\log p\), without factoring the rest of \(p-1\).

\paragraph{Proof.}
The subgroup \(G=(\mathbb F_p^\times)^r\) is proper.
If \(p\leq2(\log p)^2\), enumerate all nonzero residues.
Otherwise Bach's Theorem~2 supplies an integer outside \(G\) below that
cutoff, which is then below \(p\). In either case test candidates by
\(a^{(p-1)/r}\ne1\). Writing \(a=g^v\) in the cyclic group shows that this
is equivalent to \(r\nmid v\), exactly the desired condition.
The search has \(O((\log p)^2)\) candidates, with polynomial bit cost per
test. The analytic character with kernel \(G\) need not be computed.
\(\square\)

Bach's Theorem~2 quantifies over all positive integers below its bound.
Theorem~3, also p.~372, quantifies only over integers coprime to the modulus
and has constant \(3\). The preceding argument handles the possible zero
residue and does not conflate these statements. A fallback nonspecific
constant follows from Shoup's ordinary-character Proposition~9,
author-preprint p.~10 \cite{CXNShoup92}.

\subsection{CX-N03: certified primitive roots with factored order}
\label{cx-n03}

Given \(p\) and the complete prime factorization
\(p-1=\prod_{i=1}^t r_i^{e_i}\), a multiplicative generator of
\(\mathbb F_p^\times\) can be constructed in deterministic polynomial time
on \(\mathsf D\).

\paragraph{Proof.}
Use CX-N02 to find \(a_i\) with \(a_i^{(p-1)/r_i}\ne1\), and set
\[
              b_i=a_i^{(p-1)/r_i^{e_i}},\qquad b=\prod_i b_i .
\]
If \(a_i=g^{v_i}\), then \(r_i\nmid v_i\), so \(b_i\) has exact order
\(r_i^{e_i}\). In an abelian group, elements of pairwise coprime orders
have product of that product order: projection to every primary component
shows that any annihilating exponent is divisible by every component
order. Thus \(b\) has order \(p-1\). Since \(t\leq\log_2(p-1)\), the cost
is polynomial. The supplied factorization also certifies any candidate
through the tests \(b^{(p-1)/r_i}\ne1\). \(\square\)

No deterministic polynomial-time integer-factorization theorem is
provided here. A short guaranteed candidate list and efficient
identification of a generator are different outputs.

\subsection{CX-N04: Shoup's least primitive root bound}
\label{cx-n04}

For odd prime \(p\), Shoup's Theorem~3 gives the least positive primitive root modulo \(p\) as
\[
 O\!\left(t^4(\log t+1)^4(\log p)^2\right),\qquad t=\omega(p-1).
\]
The exact author-preprint locators are Theorem~3, p.~4; its proof in
\S6, p.~13; and Proposition~9, p.~10 \cite{CXNShoup92}.
All analytic characters used here are rational Dirichlet characters
modulo \(p\).

\paragraph{Proof of transfer and output.}
Apply \(\mathsf D\) to the finite family of character powers in the proof.
The remaining finite sieve is unconditional, so the bound and its
effective constants remain valid. The product of \(t\) different primes
is at least \(t!\); hence \(t(\log t+1)=O(\log p)\). The displayed bound is
therefore \(O((\log p)^6)\). Enumeration gives a polynomial-size list
containing a primitive root. Supplied factorization of \(p-1\) makes
recognition polynomial by CX-N03. The small bound alone does not provide
that factorization or a replacement recognition algorithm. \(\square\)

\subsection{CX-N05: Shoup's quadratic-extension candidate construction}
\label{cx-n05}

For odd prime \(p\), Shoup's different Theorem~2 constructs a candidate set for a primitive
root of \(\mathbb F_{p^2}\): author-preprint p.~3, proof \S5, pp.~12--13,
and Proposition~13 with the ray-class definitions in \S4, pp.~10--12
\cite{CXNShoup92}. Choose a prime \(\delta\equiv1\pmod4\) for which
\(-\delta\) is a quadratic nonresidue modulo \(p\), and set
\[
 K=\mathbb Q(\sqrt{-\delta}),\quad
 \mathcal O_K=\mathbb Z[\sqrt{-\delta}],\quad
 \mathfrak m=p\mathcal O_K .
\]
Proposition~10 uses rational characters to give
\(\delta=O((\log p)^2)\). Inertness identifies
\(\mathcal O_K/\mathfrak m\) with \(\mathbb F_{p^2}\).
The later character sums concern finite-order characters of
\(I_{\mathfrak m}/P_{\mathfrak m}^{+}\). Positivity is vacuous because
\(K\) is imaginary. These are Hecke characters over \(K\), together with
the trivial factor \(\zeta_K\). RH for the ray class field zeta
\(\zeta_{H_{K,\mathfrak m}}\) is sufficient.

\paragraph{Full transfer.}
Conditional \(\mathsf H\) covers every required character, so the
published character-sum and sieve estimates remain valid. The enumeration
is polynomial in \(\log p\), using the stated bounds on \(\delta\) and
the class number. Reduction gives the candidate set. An unconditional
deterministic field-isomorphism algorithm transfers it to a supplied
model \cite[Theorem~(1.2), p.~330]{CXNLenstra91}. With the factorization of \(p^2-1\), order tests identify a
generator; with only a specified prime \(r\mid p^2-1\), a power test
identifies an \(r\)-nonresidue in the list. \(\square\)

Here \(\mathsf D\) gives the initial prime selection and \(\zeta_K\),
because \(K/\mathbb Q\) is quadratic, but does not by this argument
supply the ray-class factors. The ray class field need not be abelian
over \(\mathbb Q\). This concerns the stated \emph{primitive-root candidate}
proof. It is not a necessity claim for each weaker task:
CX-N13 gives a smaller-hypothesis replacement for every specified
prime \(r\mid p-1\) in all explicitly given \(\mathbb F_{p^m}\),
including quadratic nonresidues.

\subsection{CX-N11: an explicit auxiliary Kummer field}
\label{cx-n11}

Let \(p\) be prime, \(n>1\), and \(p\nmid n\). Set
\[
              F=\mathbb Q(\zeta_{2n}),\qquad E=F(p^{1/n}).
\]
RH for \(\zeta_E\) and \(\zeta_{\mathbb Q}\) suffices for a prime
\[
 q=O(n^4\log^2(np)),\quad q\equiv1\pmod{2n},\quad
 \gcd\!\left((q-1)/\operatorname{ord}_q(p),n\right)=1.             \tag{N11.1}
\]
A search for it has polynomial cost in \(n,\log p\).
This supplies Buchmann--Shoup's Fact~2.1, p.~1314, and the closely
related Adleman--Lenstra Proposition~3, p.~351
\cite{CXNBuchmannShoup96,CXNAdlemanLenstra86}.

\paragraph{Exact analytic input.}
Bach--Sorenson's Theorem~5.1, p.~1729, gives an unramified rational prime
with any prescribed Frobenius conjugacy class in \(E/\mathbb Q\),
\(E\ne\mathbb Q\), satisfying
\[
              q\leq(4\log|\Delta_E|+2.5[E:\mathbb Q]+5)^2.      \tag{N11.2}
\]
Their convention in \S2, p.~1720 after (2.10), assumes RH for
\(\zeta_E\) and \(\zeta_{\mathbb Q}\) \cite{CXNBachSorenson96}.
There is no assumption of general Artin holomorphy.

\paragraph{Full proof of the field, conjugacy, and power conditions.}
The prime \(p\) is unramified in \(F\).
For odd \(n\), \(F=\mathbb Q(\zeta_n)\), which includes the case \(p=2\).
For even \(n\), \(p\) is odd and coprime to \(2n\).
Thus \(v_{\mathfrak P}(p)=1\) at every \(\mathfrak P\mid p\).
The polynomial \(X^n-p\) is Eisenstein over the discrete valuation ring
\(\mathcal O_{F,\mathfrak P}\), hence irreducible over \(F\).
Writing \(\alpha=p^{1/n}\), we obtain \([E:F]=n\).
The field \(E/\mathbb Q\) is Galois, since it contains all
\(\zeta_n^j\alpha\) and the \(2n\)-th roots of unity.

There is an \(F\)-automorphism \(\sigma(\alpha)=\zeta_n\alpha\) of order
\(n\). If \(\tau(\zeta_n)=\zeta_n^u\), then
\(\tau\sigma\tau^{-1}=\sigma^u\), as checked on \(F,\alpha\);
indeed \(\tau(\alpha)=\zeta_n^b\alpha\) for some \(b\).
The exponent \(u\) is a unit modulo \(n\). Thus each member of the
conjugacy class of \(\sigma\) fixes \(F\) and generates
\(\operatorname{Gal}(E/F)\).

Apply (N11.2) to this class. The resulting \(q\) splits in \(F\).
For even \(n\), whose cyclotomic conductor is \(2n\), this gives
\(q\equiv1\pmod{2n}\).
For odd \(n>1\), splitting gives \(q\equiv1\pmod n\), and \(q\ne2\)
since \(2\not\equiv1\pmod n\). As \(q\) is odd, it is again
\(1\pmod{2n}\). Also \(q\nmid pn\): the prime \(p\) ramifies in \(E/F\)
by the Eisenstein construction, and primes dividing \(n\) ramify in \(F\).

Choose a prime \(\mathfrak Q\mid q\) whose Frobenius is \(\sigma^u\).
In its residue field,
\(\alpha^q=\zeta_n^u\alpha\), so
\[
                    p^{(q-1)/n}=\alpha^{q-1}=\zeta_n^u .
\]
The reduction of \(\zeta_n\) has order \(n\), because \(q\nmid n\).
For every prime \(r\mid n\), therefore,
\[
                     p^{(q-1)/r}=(\zeta_n^u)^{n/r}\ne1.
\]
Writing \(p=g^a\) in \(\mathbb F_q^\times\), this says \(r\nmid a\).
Since \(\operatorname{ord}_q(p)=(q-1)/\gcd(a,q-1)\), it proves the gcd
condition in (N11.1).

Finally \(|\Delta_F|\leq(2n)^{[F:\mathbb Q]}\), and the relative
discriminant divides the polynomial discriminant ideal
\((n^np^{n-1})\). The discriminant tower formula gives
\[
 \begin{split}
 \log|\Delta_E|
 &\leq n[F:\mathbb Q]\log(2n)
       +[F:\mathbb Q](n\log n+(n-1)\log p)\\
 &=O(n^2\log(np)),\qquad [E:\mathbb Q]\leq2n^2.
 \end{split}
\]
Equation (N11.2) gives (N11.1), with effective constants.
Enumeration up to this bound is polynomial in \(n,\log p\);
primality and order testing, even using trial division of the
polynomially bounded \(q-1\), retain that cost. \(\square\)

The unique degree-\(n\) field \(K\subset\mathbb Q(\zeta_q)\) is real
because \(n\mid(q-1)/2\). The image of \(p\) in the cyclic quotient
\(\operatorname{Gal}(K/\mathbb Q)\) has order \(n\) by (N11.1);
hence \(p\) is inert and \(\mathcal O_K/p\mathcal O_K\cong\mathbb F_{p^n}\).
Adleman--Lenstra's Theorem~1, p.~350, Algorithm~A, pp.~350--351,
and Lemma~6, p.~352, use this mechanism to construct a field model.
That task is different from factoring an arbitrary input polynomial.

The extension \(E/F\) is cyclic. Equivalently, its finite family of
Hecke characters over \(F\) supplies \(\zeta_E\). Although \(\mathsf D\)
covers \(\zeta_F\), the additional Kummer characters are not thereby
rational Dirichlet characters. Conditional \(\mathsf H\) covers them.

\subsection{CX-N06: Buchmann--Shoup generating sets}
\label{cx-n06}

Buchmann--Shoup construct a model of \(\mathbb F_{p^n}\) and a set
generating its multiplicative group in deterministic time
\[
                              (n\log p)^{O(n)}.                \tag{N06.1}
\]
The main statement is on pp.~1311--1312; p.~1314 explicitly defines
ERH as RH for Dedekind zeta functions of all number fields
\cite{CXNBuchmannShoup96}.
For \(p\nmid n\), their auxiliary field \(K\) is the real cyclic field of
CX-N11. Fact~4.3, p.~1319, invokes Bach's Theorem~4 for the ray class
group modulo \(p\mathcal O_K\). A concrete sufficient package is
\[
                   \zeta_{\mathbb Q},\quad\zeta_E,\quad
                   \zeta_{H_{K,p\mathcal O_K}} .
\]
The ray modulus has no added infinite places, as in the paper.

\paragraph{Full transfer.}
Conditional \(\mathsf H\) supplies these zeta hypotheses. CX-N11
discharges prime selection; CX-N07 discharges ray-class generation.
The remaining lattice and finite-field operations preserve (N06.1).
For \(n=1\) and \(p>2\), Bach's Theorem~3 gives a group-generating set of size \(O((\log p)^2)\); \(p=2,n=1\) is trivial. If \(p\mid n\), then \(p\leq n\); enumerating irreducible-polynomial candidates and then all field elements costs \(p^{O(n)}\leq n^{O(n)}\), still within the stated bound. If \(S\) generates \(\mathbb F_{p^n}^\times\), then
for any specified prime \(r\mid p^n-1\), some member lies outside the
proper subgroup of \(r\)-th powers. Scanning with the exponent test
finds it. With complete factorization of \(p^n-1\), isolate primary
orders and multiply as in CX-N03. A generating set need not contain
one primitive root. \(\square\)

CX-N13 supplies a stronger nonresidue consequence for primes \(r\mid p-1\)
when the field model is already supplied.
For fixed \(n\), the time is polynomial in \(\log p\); for variable
\(n\), the displayed bound is not polynomial in the usual input
parameters. Discharging GRH does not change that degree dependence.

\subsection{CX-N07: class-group generators and their exact class field}
\label{cx-n07}

Bach's class-group bound is
\[
        \operatorname{Cl}(\mathcal O_K)=
        \langle[\mathfrak p]:N\mathfrak p\leq
                             12(\log|\Delta_K|)^2\rangle .     \tag{N07.1}
\]
The original locators are Theorem~4, p.~375, and the corollary on p.~376
\cite{CXNBach90}. Its analytic family comprises \(\zeta_K\) and the
nontrivial ideal-class-character \(L_K(s,\chi)\), which are unramified
finite-order Hecke \(L\)-functions over \(K\).

\paragraph{Proof from the precise character theorem.}
If those prime classes generate a proper subgroup \(G\), take a
nontrivial character of the quotient
\(\operatorname{Cl}(\mathcal O_K)/G\). Its conductor is \(1\).
The first bound in Bach's Theorem~4 gives a prime with
\[
 \chi(\mathfrak p)\ne1,\qquad
 N\mathfrak p\leq3\log^2(|\Delta_K|^2)=12(\log|\Delta_K|)^2,
\]
contradicting triviality on \(G\). The omitted degree-one case has
trivial class group. \(\square\)

\paragraph{Full class-field transfer.}
For the Hilbert class field \(H_K\), reciprocity gives
\[
 \operatorname{Gal}(H_K/K)\cong\operatorname{Cl}(\mathcal O_K),
 \qquad
 \zeta_{H_K}(s)=
  \prod_{\chi\in\operatorname{Cl}(\mathcal O_K)^\wedge}L_K(s,\chi).
                                                                    \tag{N07.2}
\]
The trivial factor is \(\zeta_K\); the other factors are entire.
None has a pole in \(0<\Re s<1\). Thus an off-line zero of any factor
would be an off-line zero of the product. RH for \(\zeta_{H_K}\)
suffices for all these factors. Conditional \(\mathsf H\) supplies it,
and hence (N07.1).

For a ray modulus \(\mathfrak m\), the same argument uses the ray
class field \(H_{K,\mathfrak m}\). Extra Euler factors from imprimitive
characters have zeros only on \(\Re s=0\), so do not affect the argument.
\(\square\)

Belabas--Diaz y Diaz--Friedman's refined computed bound assumes RH
only for the nontrivial class-character factors: author-preprint
Theorem~1.1 p.~3; Theorem~2.1 p.~5; proof pp.~4--6
\cite{CXNBDF08}. Their introduction distinguishes this from Bach's use
of both these factors and \(\zeta_K\). Neither gate should simply be
replaced by RH for \(\zeta_K\).

Even if \(K/\mathbb Q\) is abelian, so that \(\mathsf D\) covers
\(\zeta_K\), \(H_K/\mathbb Q\) need not be abelian. The additional
class-character gate remains in this proof. The bound gives generators,
not their relations. Algorithmic enumeration must also specify the
maximal order and its arithmetic, or account for computing them;
a defining polynomial is not silently the same input.

\subsection{CX-N08: verification of class groups and units}
\label{cx-n08}

Let \(S\) contain all archimedean places and a finite factor base.
Write \(K_S\) for the \emph{group of \(S\)-units} and \(H\subset K_S\)
for the subgroup generated by sampled \(S\)-units. For the valuation map
\[
 \phi:H\longrightarrow\mathbb Z^{S_{\rm fin}},\qquad
 x\longmapsto(\operatorname{ord}_{\mathfrak p}x)_{\mathfrak p},
\]
the tentative units and class group are \(\ker\phi\) and
\(\operatorname{coker}\phi\). In particular, \(K_S\) is not a relation
lattice. The exact source is Lenstra \S5, equation (5.7), author-preprint
p.~25 \cite{CXNLenstra92}.

\paragraph{Full verification proof.}
Assume the factor-base classes generate, \(H\) contains all roots of
unity, and the tentative groups have the required finite-index and rank
properties. The exact sequence
\[
 0\longrightarrow\ker\phi\longrightarrow\mathcal O_K^\times
 \longrightarrow K_S/H\longrightarrow\operatorname{coker}\phi
 \longrightarrow\operatorname{Cl}(\mathcal O_K)\longrightarrow0
\]
implies
\[
                              h'R'=hR[K_S:H],                 \tag{N08.1}
\]
where primes denote tentative class number and regulator.
The unit index is \(R'/R\), and the finite class-group index is
\(h'/h\); the exact sequence multiplies them to \([K_S:H]\).
If certified bounds give \(a/2<hR<a\) and the computed product also
satisfies \(a/2<h'R'<a\), then
\[
                         1\leq[K_S:H]=(h'R')/(hR)<2.
\]
The index is an integer, hence \(1\), proving completeness of both
units and class group. \(\square\)

Conditional \(\mathsf H\) supplies the factor-base bound. Where the
chosen analytic estimator assumes RH for \(\zeta_K\), it supplies
that premise also. The index test itself is algebraic and also works
with independently certified generators and analytic bounds.
It gives no rate for collecting sufficient relations. Any additional
smoothness heuristic or search-complexity assumption in a chosen
algorithm remains separate.

\subsection{CX-N12: a rigorously conditional randomized class-group algorithm}
\label{cx-n12}

Biasse--Erukulangara's modified Hafner--McCurley algorithm supplies a
concrete computation theorem \cite{CXNBiasseErukulangara}. We transfer
the input explicitly required by its algorithm boxes: a negative
\emph{fundamental} discriminant \(-d\), giving the full group structure
of \(\operatorname{Cl}(\mathcal O_{\mathbb Q(\sqrt{-d})})\).
The Main Theorem and Conjecture~1 are on author-preprint p.~4,
Algorithm~4 on p.~15, and Corollary~1 on p.~16.
Conjecture~1 explicitly concerns Hecke \(L\)-functions over number fields.

Conditional on \(\mathsf H\), the algorithm has Las Vegas expected
bit complexity
\[
 \exp\!\left(\left(\frac3{\sqrt8}+o(1)\right)
                    \sqrt{\log d\,\log\log d}\right).           \tag{N12.1}
\]
It remains randomized. The displayed bound is superpolynomial and
subexponential in the bit length \(\log d\).

\paragraph{The visible number-field family.}
The expansion step, Proposition~1, author-preprint p.~8, uses
Jao--Miller--Venkatesan's Theorem~1.1 and Corollary~1.3
\cite[arXiv version, pp.~3--4; proof pp.~9--10]{CXNJMV09}.
For modulus \(1\) in imaginary quadratic \(K\), their characters are
exactly those of \(\operatorname{Cl}(\mathcal O_K)\), and their proof
explicitly estimates the corresponding Hecke \(L\)-functions.
Thus RH for \(\zeta_{H_K}\) supplies this step and class generation
by CX-N07. The rational quadratic character of \(K\) controls the
separate quadratic splitting and class-number estimates.
The fact that \(K/\mathbb Q\) is quadratic does not itself eliminate
the class-character gate.

For completeness, the split-prime factor base in the algorithm can be
obtained from the same family. Regard each class character as an
imprimitive character modulo \(d\mathcal O_K\), extended by zero at
primes dividing \(d\). The second bound of Bach's Theorem~4 gives
\(\chi(\mathfrak p)\notin\{0,1\}\) with
\[
 N\mathfrak p\leq12\log^2(d^2N(d\mathcal O_K))
              =192(\log d)^2.
\]
Such a prime is unramified; it cannot be an inert prime ideal, since
an inert prime ideal is the principal ideal \((p)\) and has class
character value \(1\). It therefore lies over a split rational prime.
Separating a proper generated subgroup by a character proves that
these split-prime classes generate. Choosing only one of the two
primes above each split rational prime preserves generation, because
the two classes are inverses. Only finite Euler factors have
been added, so no additional number-field character family is needed.

\paragraph{Full transfer, certification, and expected time.}
Conditional \(\mathsf H\) supplies the precise hypothesis explicitly
declared in Conjecture~1. Hence the paper's conditional factor-base,
relation-probability, and class-number estimates apply with the
published parameters.

Let \(\Lambda\subset\mathbb Z^n\) be the full relation lattice for the
factor base, and \(\Lambda'\subset\Lambda\) the collected full-rank
lattice. Every collected vector is a true relation. Since the factor
base generates,
\[
       h'=\det\Lambda'=h[\Lambda:\Lambda'],\qquad
       h=|\operatorname{Cl}(\mathcal O_K)|.
\]
Algorithm~4, steps 15--20, computes a certified \(h^*\) with
\(h^*\leq h<2h^*\), obtains \(h'\) from the Smith normal form, and
declares \textsc{Failure} if \(h'\geq2h^*\).
An accepted output satisfies \(1\leq h'/h<2h^*/h\leq2\);
the integer index is therefore \(1\). Every accepted group is correct
under the analytic estimates, and unsuccessful branches are declared
failures. The Smith-form procedure is itself Las Vegas.

Put \(L(d)=\exp(\sqrt{\log d\,\log\log d})\).
Corollary~1 supplies a time-bounded run of cost
\(L(d)^{3/\sqrt8+o(1)}\) with success probability at least
\(1-1/d^{1+o(1)}\). Independent restarts on failure have a geometric
number of trials, with expectation at most \(2\) for sufficiently
large \(d\); the bounded smaller inputs can be handled directly.
This gives (N12.1) with its randomness retained. The corollary balances
\[
             \min_{z>0}\max\{3z,z+1/(4z)\}=3/\sqrt8,
             \qquad z=1/\sqrt8.
\]
The expectation deduction is justified by detectable failures and
correctness of accepted outputs, not by high success probability alone.
\(\square\)

\subsection{CX-N09: general finite-field factorization}
\label{cx-n09}

Evdokimov's published theorem gives deterministic factorization of a
degree-\(d\) polynomial over an explicit field \(\mathbb F_q\) in
\[
                            (d^{\log d}\log q)^{O(1)}            \tag{N09.1}
\]
under the relevant GRH. The primary statement and radical-oracle
description are in the abstract on p.~209 \cite{CXNEvdokimov94}.
The degree dependence is quasipolynomial. For fixed \(d\), the bound
is polynomial in \(\log q\); it is not a general polynomial bound
in \(d,\log q\). Unconditional randomized polynomial-time factorization
already predates this result \cite[introduction, author pp.~1--2]{CXNGao01}.

\paragraph{A sufficient broad radical gate.}
Evdokimov's earlier Theorem~I.3, pp.~109--110, gives deterministic
\(k\)-th root extraction in an explicitly given \(\mathbb F_{p^m}\)
in time polynomial in \(k,m,\log p\) \cite{CXNEvdokimov89}.
Its proof uses field construction and isomorphisms from
Theorems~I.1 and I.2, pp.~107--109. The analytic convention on p.~105
explicitly allows \emph{Dirichlet characters of abelian number fields}.
These are finite-order Hecke characters over those fields.
For example, case III of Theorem~I.1, p.~108, uses the quadratic
residue character modulo \(p\) over \(K=\mathbb Q(i)\), identifying
its ideal character using the fact that \(\mathbb Z[i]\) is a PID.
The hypothesis is not automatically the rational Dirichlet assertion.

Every character allowed by that oracle theorem belongs to
\(\mathsf H\). Conditional on \(\mathsf H\), the radical oracle has
the required polynomial bound. Substituting it into the published
radical-oracle reduction gives (N09.1), with its degree dependence
unchanged. This uses no assumption of arbitrary Artin holomorphy.
It is a transfer of the cited algorithmic theorem, not an independent
reproof of the complete Evdokimov algebraic algorithm.
CX-N15 supplies an independent composition of the same bound from
accessible primary algorithm statements, with explicit Kummer-zeta
hypotheses and without that original internal reduction.

\paragraph{Why a smaller-hypothesis transfer needs its own audit.}
Gao proves that radicals in an elementary algebra
\(\mathbb F_q^N\) require only \(r\)-nonresidues in the base
\(\mathbb F_q\), for primes \(r\mid\gcd(k,q-1)\):
\S2.4, author-preprint pp.~7--9, Lemma~2.6 \cite{CXNGao01}.
His introduction, pp.~1--2, records the unconditional reduction to
completely split prime-field polynomials.
These facts make a rational-character route a legitimate question;
they do not alone audit subsequent automorphism splitting.
Gao's Lemma~3.6, author-preprint p.~16, invokes the latter from
R\'onyai's work.

The broad \(\mathsf H\)-transfer is therefore established here, while
a \(\mathsf D\)-only transfer of the entire general theorem is not
asserted. The separately proved CX-N14 handles several specified
automorphism-order cases, including prime-order powers whose order divides
\(p-1\), using at most rational characters. The unrestricted step
remains a stated gate in this audit.

\subsection{CX-N10: special factorization families preserve their parameters}
\label{cx-n10}

Bach--von zur Gathen--Lenstra's Theorem~1, p.~6, is a precise
special-field theorem \cite{CXNBGL01}. Let \(S(p^k)\) be the prime
divisors of the integer \(\Phi_k(p)\), \(s(p^k)\) its largest prime
divisor, and \(R(p^k)\) the primes dividing \(\ell-1\) for some
\(\ell\in S(p^k)\). For an explicit \(\mathbb F_{p^n}\) and \(f\),
the theorem assumes a degree-\(r\) irreducible polynomial over
\(\mathbb F_p\) for each \(r\in R(p^k)\) not dividing \(n\), and has time
\[
                     (s(p^k)+\deg f+n\log p)^{O(1)}.            \tag{N10.1}
\]

\paragraph{Transfer.}
Adleman--Lenstra constructs the extra irreducible-polynomial inputs
when its auxiliary-prime hypothesis is supplied.
For each required prime \(r\ne p\), CX-N11 supplies the concrete
field \(\mathbb Q(\zeta_{2r},p^{1/r})\); characteristic-\(p\)
extensions use the separate algebraic construction.
Conditional \(\mathsf H\) covers these hypotheses and removes the
extra input gates while retaining (N10.1).
The same paper's Theorem~2, p.~7, constructs a primitive
\(\ell\)-th root of unity from specified irreducible-polynomial inputs;
Theorem~3, p.~7, constructs the relevant nonresidues from roots of
unity. These intermediate tasks have explicit inputs and parameters.
They do not give a polynomial bound for arbitrary polynomial
factorization. \(\square\)

For Fermat or Mersenne characteristics, \(k=1\) or \(k=2\) makes
\(\Phi_k(p)\) a power of \(2\), leaving \(R(p^k)\) empty.
The corollary on p.~6 is then unconditional, rather than newly
unconditional by either GRH claim.

R\'onyai's other results also have structural parameters.
The primary abstract of \cite{CXNRonyai92} gives time polynomial in
the input degree, coefficient size, \(\log p\), and the rational
splitting-field degree; the latter may be exponential in the usual
input length. Gao's introduction, p.~2, records R\'onyai's 1988
algorithm for finding a nontrivial factor of a degree-\(d\) completely
split polynomial, with cost polynomial in \(d^r,\log p\), for a
prime divisor \(r\mid d\). It is polynomial for bounded \(r\), not
uniformly when \(r\) grows. These are sourced parameter distinctions.
The full original proofs and precise theorem numbers of these two
R\'onyai statements were not independently retrieved; no stronger
rational-Dirichlet transfer is inferred from the acronym GRH.


% CX_GRH supplement, for the number-algorithm component.
% Finding ID CX-N13 agreed with /root/number_algorithms.
% Cite BIMS17 from koiran_bibliography.tex.

\subsection{CX-N13: nonresidues for primes dividing \texorpdfstring{$p-1$}{p-1} in supplied extensions}
\label{cx:n13}

\paragraph{Statement.}
Let $p$ be a prime, let $r$ be a specified prime divisor of $p-1$, and let
$F=\mathbb F_p[X]/(f)$ be supplied by an irreducible polynomial $f$
of degree $m$. An $r$-th nonresidue in $F$ can be constructed in
deterministic $\operatorname{poly}(m\log p)$ time under the stipulated
Dirichlet-GRH theorem $H_A$. In particular this holds for quadratic
nonresidues in every supplied extension of every odd prime field.
For $r=2$ and even $m$, the construction is already unconditional.

\paragraph{Established even-degree input.}
Bhargava--Ivanyos--Mittal--Saxena \cite[Corollary~1.3, PDF p.~2]{BIMS17}
gives deterministic $\operatorname{poly}(m\log p)$-time $r$-th-root
extraction when $r$ is constant and $r\mid\gcd(m,p-1)$.
The following remark and \S2 specify that the irreducible polynomial
defining the field is part of the input. For $r=2$, this covers every
even $m$ and every odd $p$. This is a prior unconditional theorem,
not a consequence newly obtained from either RH premise.
For variable $r\mid p-1$, their Theorem~1.2 also gives the transfer
once $\zeta_r$ is available and $r\mid m$: the supplied irreducible
$f$ satisfies their Property~1.1. We give the direct algebraic
argument for odd $r$, so the uniformity does not rely on a
constant-$r$ root-of-unity algorithm.

\begin{proof}[Proof of the all-degree transfer and output refinement]
Under $H_A$, CX-N02 supplies $a\in\mathbb F_p^\times$ with
$a^{(p-1)/r}\ne1$. Put $\zeta=a^{(p-1)/r}$, which has order $r$.
If $r\nmid m$, then, since $p\equiv1\pmod r$,
\[
 a^{(p^m-1)/r}
 =\zeta^{1+p+\cdots+p^{m-1}}
 =\zeta^m\ne1.
\]
Thus the constant $a$ is the required nonresidue in the supplied field.

Suppose next that $r\mid m$ and $r$ is odd. Set $Q=p^{m/r}$ and let
$\sigma:F\longrightarrow F$ be $u\longmapsto u^Q$.
The cyclic group $F^\times$ has order $(Q-1)S$, where
$S=1+Q+\cdots+Q^{r-1}$. The image of $u\mapsto u^{Q-1}$
has order $S$. Since $Q\equiv1\pmod r$, we have $r\mid S$,
so that image contains the element $\zeta$ of order $r$.
Consequently the $\zeta$-eigenspace of $\sigma$ is nonzero.
Compute a nonzero element $u$ of this eigenspace by solving
\[
 (\sigma-\zeta I)u=0
\]
as a linear system over $\mathbb F_p$ in the supplied polynomial
basis of $F$. The Frobenius matrix and a nonzero kernel vector can
be computed in polynomial bit complexity in $m\log p$.
Then $u^{Q-1}=\zeta$. Write $Q=1+rt$. Modulo $r^2$,
\[
 1+Q+\cdots+Q^{r-1}
 \equiv r+rt\frac{r(r-1)}2\equiv r\pmod{r^2},
\]
where oddness of $r$ is used in the second congruence. Therefore
\[
 u^{(p^m-1)/r}
 =\left(u^{Q-1}\right)^{(1+Q+\cdots+Q^{r-1})/r}
 =\zeta\ne1.
\]
This constructs the required $r$-th nonresidue. In this branch
$r\le m$; in the preceding branch CX-N02 and modular powering
are polynomial in $\log p$ independently of the numerical value
of $r$. The claimed uniform polynomial bound follows.

It remains to cover $r=2$ and even $m$, and here no RH assumption
is needed. Invoke the cited unconditional square-root algorithm.
Start with $z=-1$. While $z^{(p^m-1)/2}=1$, replace $z$ by any
square root supplied by that algorithm. Initially $z$ has order $2$.
A square root of an element of order $2^j$, for $j\ge1$, has order
$2^{j+1}$: squaring halves that order, and an element with square
of $2$-power order has no odd order component. Consequently each
iteration doubles the order. After at most
$v_2(p^m-1)-1\le\log_2(p^m)$ iterations, $z$ has the full
$2$-primary order of $F^\times$ and is a nonsquare. Euler tests and
the polynomially many square-root calls retain polynomial bit
complexity in $m\log p$. This branch does not call the initial
$H_A$-dependent prime-field search. The cases exhaust all $m$.
\end{proof}

This scope correction applies to the corresponding nonresidue
consequences of CX-N05 and CX-N06. It does not remove the
number-field Hecke hypothesis from Shoup's broader primitive-root
candidate-set theorem or supply the same reduction when
$r\mid p^m-1$ but $r\nmid p-1$. The extension model is supplied;
the statement does not claim an unconditional construction of that
model from $p$ and $m$ alone. In characteristic two every element
is a square, so a quadratic nonresidue does not exist.

An independent elementary verification of the even-degree quadratic
case, including the $p\equiv3\pmod4$ trace construction, is recorded
in the working proof notes. It is not presented here as a new
discovery of the established theorem.


\subsection{CX-N14: proved cases of splitting a prime-field algebra from an automorphism}
\label{finding:CX-N14}

\paragraph{Finding.}
Let \(A\) be a commutative unital \(\mathbb F_p\)-algebra
given by an explicit basis and multiplication table, promised to
be isomorphic to \(\mathbb F_p^d\).
Let \(\tau\ne1\) be a supplied \(\mathbb F_p\)-algebra automorphism,
given by its matrix, and let \(h=\operatorname{ord}(\tau)\).
If at least one of the following conditions holds, where
\(r\) in (ii) and (iii) is a prime divisor of \(h\),
\[
\begin{array}{ll}
\textnormal{(i)} & p\leq d,\\
\textnormal{(ii)} & r\mid(p-1),\\
\textnormal{(iii)} & r\ne p,\quad
 v_r\!\left(p^{s}-1\right)=1,\quad s=\operatorname{ord}_r(p),
\end{array}
\]
then a nonzero zero divisor of \(A\) can be found deterministically
in \(\operatorname{poly}(d,\log p)\) bit operations.
Cases (i) and (iii) are unconditional.
Case (ii) uses only \(A_D\), through CX-N02.
For \(A=\mathbb F_p[X]/(f)\), the output yields a proper factor of
\(f\).

This is a consequence of the stated constructions and the precise
algebraic algorithms of
\cite[Lemma~4.2 and Proposition~4.6]{CXNIKRS08};
it is not presented as a new discovery of those source algorithms.
The finding identifies concrete parts of the general
automorphism-splitting step whose analytic requirements are
discharged or absent.

\paragraph{Order recovery and reduction to prime order.}
An automorphism of \(\mathbb F_p^d\) permutes the primitive
idempotents, so its order divides
\(T=\operatorname{lcm}(1,\ldots,d)\).
The bit length of \(T\) is polynomial in \(d\), and its prime factors
are all at most \(d\).
Starting with \(T\), divide out a prime factor whenever the
corresponding power of \(\tau\) is the identity.
Matrix powering performs these tests and determines \(h\) in
polynomial time.
The same small-prime factorization finds the prime divisors \(r\)
of \(h\).
Replace \(\tau\) by \(\tau^{h/r}\) for a chosen \(r\).
This has exact order \(r\), and \(r\leq d\).
The integer \(s=\operatorname{ord}_r(p)\), when \(r\ne p\), and
the valuation test in (iii) are also computable in polynomial
time: \(s\leq r-1\), and \(p^s\) has \(O(d\log p)\) bits.
Thus selecting a covered case requires no unprovided integer
factorization oracle.

\paragraph{Proof of case (i).}
Here \(p\leq d\); the supplied automorphism is not needed for
this case. In particular, it covers every characteristic-order
case \(r=p\), since a prime divisor of \(h\) is at most \(d\).
Choose a basis element \(b\) outside the one-dimensional scalar
subspace \(\mathbb F_p1\).
Under the promised product decomposition its coordinates are not
all equal.
For some \(a\in\mathbb F_p\), the element \(b-a1\) has a zero
coordinate, and it is nonzero for every \(a\).
Enumerate the \(p\leq d\) scalars and test the multiplication
matrices of \(b-a1\) for singularity.
A singular matrix gives a nonzero zero divisor.
This requires polynomial time and no analytic assumption.
\hfill\(\square\)

\paragraph{Proof of case (ii).}
Use CX-N02 to find an \(r\)-th nonresidue \(a\in\mathbb F_p^*\).
Set \(\zeta=a^{(p-1)/r}\), of exact order \(r\).
The prime-order permutation \(\tau\) has an \(r\)-cycle, so its
\(\zeta\)-eigenspace is nonzero.
Linear algebra finds
\[
 0\ne z\in A,\qquad \tau(z)=\zeta z.
\]
If multiplication by \(z\) is singular, it is already the desired
zero divisor. Otherwise \(z\) is a unit.

Compute \(B=A^\tau=\ker(\tau-1)\) and its multiplication table by
linear algebra.
It is a split \(\mathbb F_p\)-algebra, with components indexed by
the orbits of \(\tau\).
The unit \(c=z^r\) belongs to \(B\).
On each orbit its constant coordinate is the \(r\)-th power of
some coordinate of \(z\), so \(c\) has an \(r\)-th root in \(B\).

The elementary-algebra root algorithm in
\cite[\S2.4, Lemma~2.6, author pp.7--9]{CXNGao01}
finds such a \(y\in B\) in
\(\operatorname{poly}(d,r,\log p)\) time given a primitive
\(r^e\)-th root of unity, where \(p-1=r^e w\), \(r\nmid w\).
The same nonresidue \(a\) supplies that root as \(a^w\).
Repeated division by \(r\) computes \(e,w\).
This uses no extension-field nonresidue and no factorization of
the complementary factor \(w\).

Now
\[
 \prod_{j=0}^{r-1}(z-\zeta^j y)=z^r-y^r=0.
\]
Every factor is nonzero: \(z=\zeta^j y\) would imply, on applying
\(\tau\), that \(\zeta z=z\), contradicting invertibility of \(z\).
At least one factor is a nonunit, since a product of units cannot
be zero.
Testing the \(r\leq d\) multiplication matrices finds a nonzero
nonunit, hence a nonzero zero divisor.
All operations have polynomial bit complexity.
\hfill\(\square\)

\paragraph{Constructive removal of fixed components.}
For case (iii), and also as an unconditional preliminary test in
any prime-order case, form
\[
 J=\bigl(a-\tau(a):a\in A\bigr)\subseteq A.
\]
For an explicit basis \(b_1,\ldots,b_d\), this ideal is the linear
span of the \(d^2\) products
\(b_j(b_i-\tau(b_i))\), so it is computable by linear algebra.
It is nonzero because \(\tau\ne1\).
On a component fixed by \(\tau\), every generator vanishes.
On a moved component, an abstract primitive idempotent supported
there proves that some generator has nonzero coordinate.
Consequently \(J\) is exactly the ideal supported on the moved
components.

If \(J\ne A\), any nonzero element of \(J\) is a zero divisor,
annihilated by the nonzero complementary component ideal.
This immediately solves the problem.
If \(J=A\), every cycle of \(\tau\) has length \(r\).
The cyclic action is then semiregular in the sense of IKRS:
no nonidentity power acts as the identity on a nonzero ideal.
This establishes the required promise for the next step
constructively; no component decomposition has been assumed as
algorithmic input.

\paragraph{Proof of case (iii).}
We may assume the preceding test gave \(J=A\).
Put
\[
 D_r=\mathbb F_p[T]/(\Phi_r(T)),\qquad
 C=A\otimes_{\mathbb F_p}D_r,\qquad
 \zeta=T\bmod\Phi_r(T).
\]
The algebra \(D_r\) is semisimple because \(p\ne r\).
Every primitive \(r\)-th root has Frobenius orbit length
\(s=\operatorname{ord}_r(p)\).
It follows that \(D_r\) is a product of copies of
\(\mathbb F_{p^s}\), and so is \(C\), since \(A\) is split over
\(\mathbb F_p\).
No explicit factorization of \(\Phi_r\) is needed for this argument
or for the construction of \(C\).

The valuation-one hypothesis implies that every element of
\(r\)-power order in \(C^\times\) has order dividing \(r\).
For \(1\leq b<r\), let \(\rho_b\) fix \(A\) and send
\(\zeta\) to \(\zeta^b\).
In this exponent-\(r\) case, the Teichm\"uller subgroup used in IKRS
is exactly
\[
 \mathcal T_{A,r}
 =\{u\in C^\times:u^r=1,\ 
          \rho_b(u)=u^b\text{ for all }1\leq b<r\};
\]
their Teichm\"uller exponent lift is congruent to \(b\) modulo \(r\).

IKRS Proposition~4.6, with the semiregular action just verified,
constructs \(x\in\mathcal T_{A,r}\) satisfying
\(\tau(x)=\zeta x\) in polynomial time.
The element \(\zeta\) also belongs to \(\mathcal T_{A,r}\).
Both have order \(r\).
Moreover \(x\notin\langle\zeta\rangle\), since \(\tau\) fixes
\(\zeta\) but does not fix \(x\).
Thus \(\langle x,\zeta\rangle\cong C_r\times C_r\) is noncyclic.
IKRS Lemma~4.2 now produces a nonzero zero divisor in
\(A\), in \(\operatorname{poly}(r,\log|A|)\) time.
Its conclusion is in \(A\) itself: the source proof makes an
annihilator idempotent invariant under every \(\rho_b\), and
\(C^{\{\rho_b\}}=A\).
Since \(r\leq d\) and \(\log|A|=d\log p\), the required bit
complexity follows, without GRH.
\hfill\(\square\)

\paragraph{The remaining branch of this argument.}
Under \(A_D\), if none of the three cases applies, every relevant
prime-order branch has
\[
 p>d,\qquad
 r\nmid p(p-1),\qquad
 s=\operatorname{ord}_r(p)>1,\qquad
 r^2\mid p^s-1.
\]
The fixed-component test may still split individual instances.
For a semiregular instance in this branch, the above proof supplies
the element \(x\), but does not supply a coefficient
Teichm\"uller unit of sufficient \(r\)-power order in \(D_r\),
or an extension-field radical algorithm that replaces it.
Gao's general automorphism-splitting Lemma~3.6 cites
R\'onyai's 1992 result
\cite[Lemma~3.6, author p.16]{CXNGao01};
the exact additional analytic input in a chosen general route
must still be tracked.
The algebra \(A[\zeta_r]\) in IKRS has extension-field components
when \(s>1\).
These statements delimit this proved algorithm.
They do not assert that a Dirichlet-only alternative for the
remaining branch is impossible.

\paragraph{A precise sufficient algebraic input for the remaining branch.}
Let \(e=v_r(p^s-1)\), and suppose an element
\(\eta\in D_r^\times\), in the explicit power basis of \(D_r\),
is additionally supplied with
\[
 \operatorname{ord}(\eta)=r^e,\qquad
 \rho_b(\eta)=\eta^{\omega_e(b)}
 \quad(1\leq b<r),
 \qquad
 \omega_e(b)=b^{r^{e-1}}\bmod r^e.
\]
These promises can be checked by polynomially many algebra
operations; in particular, exact order is checked with the two
powers \(r^e\) and \(r^{e-1}\).
They say that \(\eta\) belongs to the coefficient
Teichm\"uller subgroup, with maximum possible \(r\)-power order.
Embedded in \(C\), it also belongs to \(\mathcal T_{A,r}\)
and is fixed by \(\tau\).
Every \(r\)-element of \(C^\times\), including the element \(x\)
from Proposition~4.6, has order dividing \(r^e\), because each
field component has \(p^s\) elements.
If \(\langle x,\eta\rangle\) were cyclic, its order would be
\(r^e\), and its element \(\eta\) of that order would generate
it. This would put \(x\) in \(\langle\eta\rangle\), contradicting
\(\tau(x)=\zeta x\ne x\).
Thus the group is noncyclic, and Lemma~4.2 gives a nonzero zero
divisor in \(A\) in polynomial time.
This explicitly delimits a sufficient additional input; the
construction of such an \(\eta\) in the remaining branch is not
claimed here under \(A_D\).
\hfill\(\square\)

\paragraph{A limit of the valuation-one argument.}
The valuation hypothesis in CX-N14 cannot simply be deleted from
its particular noncyclicity proof.
For an explicit example, take \(p=17\), \(r=3\), and
\(A=\mathbb F_{17}^3\).
The polynomial \(T^2+T+1\) has no root over \(\mathbb F_{17}\),
since a root would have order three in its multiplicative group
of order sixteen. Thus
\[
 D=\mathbb F_{17}[T]/(T^2+T+1)\cong\mathbb F_{17^2},
 \qquad \zeta=T\bmod(T^2+T+1).
\]
Here \(s=\operatorname{ord}_3(17)=2\) and
\(v_3(17^2-1)=v_3(288)=2\).
The cyclic group \(D^\times\) contains an element \(\eta\) of
order nine with \(\eta^3=\zeta\): start with a generator of its
order-nine subgroup and replace it by its inverse if necessary.
Set
\[
 x=(\eta,\zeta\eta,\zeta^2\eta)\in
 A\otimes_{\mathbb F_{17}}D=D^3,
 \qquad
 \tau(a_0,a_1,a_2)=(a_1,a_2,a_0).
\]
Then \(\tau(x)=\zeta x\), \(x\) has order nine, and
\(x^3=(\zeta,\zeta,\zeta)\).
The nonidentity cyclotomic automorphism \(\rho_2\) of \(D\)
is its seventeenth-power Frobenius.
On every element of order dividing nine, this acts by exponent
\(17\equiv-1\pmod9\), which is also
\(\omega_2(2)=2^3\equiv-1\pmod9\).
It follows coordinate by coordinate that
\(\rho_2(x)=x^{\omega_2(2)}\); hence \(x\) is in the required
Teichm\"uller subgroup.
Nevertheless \(\langle x,\zeta\rangle=\langle x\rangle\)
is cyclic of order nine.
This example satisfies the eigenunit properties used from
Proposition~4.6 but does not force that pair to be noncyclic.
It makes no assertion about which valid eigenunit a particular
implementation returns, or about the impossibility of a different
Dirichlet-only algorithm.
\hfill\(\square\)

\paragraph{Recovering a polynomial factor.}
If \(A=\mathbb F_p[X]/(f)\), represent the resulting nonzero zero
divisor by \(g\) of degree less than \(\deg f\).
Then \(\gcd(f,g)\) has positive degree because multiplication by
\(g\) is singular, and has degree less than \(\deg f\) because
\(g\ne0\) in \(A\). The Euclidean algorithm therefore returns
a proper factor in polynomial time.

\paragraph{Geometric restatement.}
The algebra \(A\) is the ring of functions on a constant finite
\'etale scheme of degree \(d\).
Its automorphism permutes the components.
The ideal \(J\) detects the moved-component open-and-closed
subscheme; the fixed algebra \(A^\tau\) is the orbit quotient.
In case (ii), a prime-field eigenfunction and an invariant radical
yield a proper component decomposition.
In case (iii), the cyclotomic coefficient extension admits the
stated Teichm\"uller construction, and the source descent returns
an idempotent on the original finite \'etale scheme.
The explicit algebra encoding and divisibility/valuation conditions
remain part of this geometric statement.


\subsection{CX-N15: a complete accessible composition for general factorization}
\label{finding:CX-N15}

\paragraph{Finding and source interface.}
Assume \(\mathsf H\), the all-number-field Hecke claim.
Given an explicit finite field \(\mathbb F_q\), \(q=p^m\), and
\(0\ne f\in\mathbb F_q[X]\) of degree \(d\geq2\), there is a
deterministic algorithm that returns its complete factorization in
\[
 \bigl(d^{\log d}\log q\bigr)^{O(1)}
\]
bit operations.
The composition below uses Adleman--Lenstra Algorithm~A and
Theorem~1, with its auxiliary-prime step supplied by CX-N11;
Bach--von zur Gathen--Lenstra Theorem~2; and the IKRS Main Theorem,
together with the explicitly proved coefficient-unit step of
CX-N14
\cite{CXNAdlemanLenstra86,CXNBGL01,CXNIKRS08}.
It does not require the unretrieved internal reduction of
Evdokimov's 1994 paper, or a general finite-field factorization
oracle.
This gives an independent proof of the broad transfer in CX-N09;
it does not change that theorem's degree dependence.

\paragraph{Step 1: construct a maximal coefficient unit.}
Consider a split algebra \(A_0\cong\mathbb F_p^{d_0}\) with a
supplied nonidentity automorphism and \(p>d\geq d_0\).
Reduce its automorphism to prime order \(r\leq d_0\) as in
CX-N14. The branch \(p\leq d\) is handled unconditionally in
Step~3. The following construction works for every such \(r\),
including \(r\mid p-1\). Put
\[
 s=\operatorname{ord}_r(p),\qquad Q=p^s,
 \qquad e=v_r(Q-1),\qquad R=r^e.
\]
The quantities \(s\leq r-1\), \(e\), and \(R\) are computable
in polynomial time in \(r,\log p\); no factorization of the
complementary part of \(Q-1\) is used.
Adleman--Lenstra Algorithm~A, with CX-N11, constructs an explicit
field
\(U=\mathbb F_{p^{sr}}\) in polynomial time in \(r,\log p\).
Indeed \(sr\leq r(r-1)\) and \(p\nmid sr\), because both
\(r\) and \(s\) are smaller than \(p\).
Compute the subfield
\[
 B=\ker\bigl(U\longrightarrow U,\ u\longmapsto u^Q-u\bigr)
   \cong\mathbb F_Q
\]
by \(\mathbb F_p\)-linear algebra, and compute its multiplication
table and inclusion in \(U\).
This is exactly the explicit-data convention of BGL, adopted
from Lenstra: basis structure constants suffice, and Frobenius
kernels provide explicit subfields
\cite[p.~329 and \S2, p.~332]{CXNLenstra91}.

BGL Theorem~2, p.~7, constructs a primitive \(r\)-th root
\(\zeta\in B\), given additionally an irreducible polynomial
over \(\mathbb F_p\) of degree \(\ell\) for each prime
\(\ell\mid r-1\) with \(\ell\nmid s\).
All these primes are smaller than \(p\) and at most \(r-1\).
Trial division factors \(r-1\) in polynomial time in \(r\),
and Adleman--Lenstra with CX-N11 constructs every required
auxiliary polynomial in polynomial time in \(r,\log p\).
Thus all BGL inputs are supplied, and its output and running-time
conditions apply. When \(r=2\), simply take \(\zeta=-1\);
there are no auxiliary primes.

Find a nonzero solution of
\[
 u^Q=\zeta u\qquad (u\in U)
\]
by \(\mathbb F_p\)-linear algebra.
Such a solution exists: the image of the power map
\(U^\times\to U^\times\), \(v\mapsto v^{Q-1}\), is the
cyclic subgroup of order
\[
 S=\frac{Q^r-1}{Q-1}=1+Q+\cdots+Q^{r-1}.
\]
Since \(Q\equiv1\pmod r\), the integer \(S\) is divisible
by \(r\); hence this subgroup contains \(\zeta\).
Let \(c=u^r\). Then
\[
 c^Q=(\zeta u)^r=c,
 \qquad c^{(Q-1)/r}=u^{Q-1}=\zeta\ne1.
\]
Thus \(c\in B\) is an \(r\)-th nonresidue.
Because \(B^\times\) is cyclic of order \(Q-1=R w_0\),
\(r\nmid w_0\), the element
\[
 \eta_0=c^{w_0}\in B^\times
\]
has exact order \(R\).
To verify the last assertion, write \(c=g^a\) for an abstract
generator \(g\) of \(B^\times\).
The nonresidue test says \(r\nmid a\), so
\(\operatorname{ord}(g^{a w_0})=R\).
This argument uses a generator only to prove the order; none is
an algorithmic input.

\paragraph{Step 2: descend the unit by equivariant interpolation.}
Let
\[
 D_r=\mathbb F_p[T]/(\Phi_r(T)),\qquad
 G=(\mathbb Z/r\mathbb Z)^\times,
 \qquad \rho_b(T)=T^b\quad(b\in G),
\]
and put \(\omega(a)=a^{r^{e-1}}\bmod R\).
This is the multiplicative Teichm\"uller lift of \(G\); its
values have order dividing \(r-1\) and reduce to \(a\) modulo
\(r\).
Since \(p^s\equiv1\pmod R\) and \(s\mid r-1\),
\(p\bmod R\) has order \(s\) and equals
\(\omega(p\bmod r)\).
Indeed reduction modulo \(r\) is injective on the
prime-to-\(r\) torsion of \((\mathbb Z/R\mathbb Z)^\times\).

Interpolate the unique polynomial \(P\in B[T]\),
\(\deg P<r-1\), with values
\[
 P(\zeta^a)=\eta_0^{\omega(a)}\qquad(a\in G).
\]
There are only \(r-1\) distinct nodes, so interpolation has
polynomial cost.
Write \(P^{(p)}\) for the polynomial obtained by raising each
coefficient to the \(p\)-th power.
For every \(a\in G\), the Frobenius identity and the preceding
Teichm\"uller congruence give
\[
 P^{(p)}(\zeta^{pa})
   =P(\zeta^a)^p
   =\eta_0^{p\omega(a)}
   =\eta_0^{\omega(pa)}
   =P(\zeta^{pa}).
\]
Both polynomials have degree below \(r-1\) and agree on all
\(r-1\) nodes. Hence \(P^{(p)}=P\), and
\(P\in\mathbb F_p[T]\).
These coefficients are extracted on the known prime-field
subspace of \(B\).

Set \(\eta=P(T)\bmod\Phi_r(T)\in D_r\).
Evaluation at all roots \(\zeta^a\) is injective after extending
coefficients to \(B\), because \(\Phi_r\) is separable and
splits there.
Every evaluated value is a unit of order \(R\), so
\(\eta^R=1\) and \(\operatorname{ord}(\eta)=R\), the latter
already following from evaluation at \(\zeta\).
For \(a,b\in G\),
\[
 P(\zeta^{ab})
 =\eta_0^{\omega(a)\omega(b)}
 =P(\zeta^a)^{\omega(b)}.
\]
Injectivity of evaluation therefore proves
\(\rho_b(\eta)=\eta^{\omega(b)}\).
Thus \(\eta\) is the maximum-order coefficient Teichm\"uller
unit required by CX-N14, constructed in
\(\operatorname{poly}(d_0,\log p)\) time under \(\mathsf H\).
No factorization of \(\Phi_r\) or general polynomial-factorization
oracle has been used.

For completeness, the last step of CX-N14 applies here for
every \(r\ne p\), without its prime-field nonresidue shortcut.
Its moved-component ideal either already splits \(A_0\), or
establishes a semiregular prime-order automorphism \(\tau\).
IKRS Proposition~4.6 then supplies
\(x\in\mathcal T_{A_0,r}\) with
\(\tau(x)=z x\), where \(z=T\bmod\Phi_r(T)\).
The algebra \(A_0\otimes D_r\) is a product of copies of
\(\mathbb F_Q\), so all of its units of \(r\)-power order have
order dividing \(R\).
The coefficient unit \(\eta\) has order \(R\) and is fixed
by \(\tau\). If \(\langle x,\eta\rangle\) were cyclic,
\(\eta\) would generate it, contradicting
\(\tau(x)=zx\ne x\).
IKRS Lemma~4.2 therefore returns a zero divisor in \(A_0\)
in polynomial time.
This proves the claimed splitting step for all prime orders in
the branch \(p>d\), including \(s=1\) and \(r=2\), without
adding an unstated Dirichlet-character premise.

\paragraph{Step 3: from a general polynomial to a split algebra.}
Polynomial gcds and derivatives give squarefree preprocessing.
If a positive-degree polynomial has zero derivative, it is a
\(p\)-th power: take the coefficient \(p\)-th roots using the
inverse Frobenius \(a\mapsto a^{p^{m-1}}\), divide the exponents
by \(p\), and retain the multiplicity \(p\).
Otherwise a nontrivial gcd with the derivative yields smaller
factors to process recursively.
These operations have polynomial bit complexity.
It therefore suffices to factor a monic squarefree polynomial
\(f\) of degree at most \(d\).

Compute the multiplication table of
\(A=\mathbb F_q[X]/(f)\), viewed as an
\(\mathbb F_p\)-algebra of dimension \(m\deg f\), and then
compute
\[
 B_f=\ker\bigl(A\longrightarrow A,\ a\longmapsto a^p-a\bigr).
\]
The map is \(\mathbb F_p\)-linear, so its kernel and the kernel's
multiplication table are obtained by linear algebra.
Abstractly, if \(f=\prod_{i=1}^t f_i\) is its distinct
irreducible factorization, then
\[
 A\cong\prod_{i=1}^t\mathbb F_{q^{\deg f_i}},
 \qquad B_f\cong\mathbb F_p^t.
\]
This proves the split promise without knowing the factors.
The primitive idempotents of \(B_f\) and \(A\) coincide.
If \(\dim_{\mathbb F_p}B_f=1\), the polynomial is irreducible.

If \(p\leq d\), choose \(b\in B_f\setminus\mathbb F_p1\)
and enumerate \(b-a1\) for \(a\in\mathbb F_p\).
Each is nonzero, and at least one has a zero component in
\(B_f\cong\mathbb F_p^t\).
Testing multiplication matrices therefore finds a nonzero zero
divisor after at most \(p\leq d\) tests.
This early branch works even when \(t<p\).

For \(p>d\), apply the IKRS Main Theorem, p.~3, to the explicit
\(t\)-dimensional algebra \(B_f\).
It deterministically gives a decomposition into ideals, each
with an automorphism whose order equals that ideal's dimension,
in time \(\operatorname{poly}(t^{\log t},\log p)\).
If the decomposition has more than one summand, the identity of
any proper nonzero summand is a nonzero zero divisor of \(B_f\);
that identity is found by linear algebra in the supplied ideal.
If there is only one summand, its automorphism has order
\(t>1\), and the preceding steps with CX-N14 produce a nonzero
zero divisor of \(B_f\) in polynomial additional time.
Thus in either case we obtain a nonzero zero divisor \(b\in A\),
and \(\gcd(f,b)\), using its polynomial representative, is a
proper nonconstant factor.
Equivalently, a primitive idempotent \(e_i\) identifies the
component factor as \(\gcd(f,1-e_i)\).
Recurse until all factors are irreducible, and reinsert the
recorded multiplicities.

\paragraph{Running time and exact analytic scope.}
There are at most \(2d-1\) nodes in the nontrivial factorization
tree, and every split algebra has dimension at most \(d\).
Linear algebra in the original quotient has size polynomial in
\(d,m,\log p\), hence in \(d,\log q\).
At each node the coefficient construction has polynomial cost;
the IKRS call has cost
\(d^{O(\log d)}(\log q)^{O(1)}\).
Multiplying by the number of nodes preserves the displayed
\((d^{\log d}\log q)^{O(1)}\) bound.

The analytic assumptions enter the constructed fields and
auxiliary irreducible polynomials through Adleman--Lenstra's
auxiliary-prime step only.
An explicit sufficient family is RH for \(\zeta_{\mathbb Q}\)
and for
\[
 \zeta_{E_\nu},\qquad
 E_\nu=\mathbb Q(\zeta_{2\nu},p^{1/\nu}),
\]
where \(\nu=sr\) or \(\nu=\ell\) for one of the auxiliary
primes \(\ell\mid r-1\), \(\ell\nmid s\), used at a node.
CX-N11 supplies the required auxiliary prime from precisely
these Dedekind-zeta hypotheses, all implied by \(\mathsf H\).
Every \(\nu\leq d^2\) is coprime to \(p\), because these
constructions occur only after the branch \(p\leq d\) has
already been handled, and the factors of \(sr\) are smaller
than \(p\).
This is a sufficient family, not a necessity assertion.
BGL's root-of-unity algorithm with those inputs, the
Frobenius/eigenvector calculations, interpolation, and IKRS
are unconditional algebraic steps.
The construction therefore supplies the formerly isolated
coefficient-unit gate under \(\mathsf H\) without circularly
calling general factorization or assuming a polynomial bound
for its varying degree.
\hfill\(\square\)

\paragraph{Finite-\'etale restatement.}
For squarefree \(f\), \(\operatorname{Spec}A\) is a finite
\'etale scheme over \(\mathbb F_p\).
Its Frobenius-fixed function algebra \(B_f\) is the constant
finite \'etale algebra indexing its connected components.
Idempotents found in \(B_f\) pull back to the corresponding
open-and-closed decomposition of \(\operatorname{Spec}A\).
The proof constructs these component projectors with the stated
quasipolynomial dependence on the input polynomial degree.


% Insert near CX-N05/N06/N11; this is a source-gate refinement, not a new finding ID.
\paragraph{The auxiliary analytic premise in the AL86/BGL route.}
Adleman--Lenstra's Proposition~3, p.~351 \cite{CXNAdlemanLenstra86},
uses a variant of Bach--Shallit's FOCS~1985 Theorem~2 to find the
small auxiliary prime needed for irreducible-polynomial construction.
In the full 1989 paper \cite[\S8, Lemma~13, pp.~213--214]{CXBachShallit89},
the analytic statement explicitly assumes RH for the Dedekind zeta
function of the top number field $L'$ in an abelian extension
$L'/L$. It supplies a degree-one prime of $L$ with prescribed Artin
symbol and norm $O(\log^2|\operatorname{disc}(L')|)$; Lemma~14 and
Theorem~3, p.~214, then bound the auxiliary rational prime.
Being abelian over $L$ does not make $L'$ abelian over $\mathbb Q$.
Thus this primary-source chain has a verified $H_M$ gate, also
realized by the explicit Kummer construction in CX-N11.

BIMS17's \cite[Conjecture~6.3 and Theorem~6.7, pp.~7--8]{BIMS17}
weak-GRH argument expressly concerns rational Dirichlet characters
and proves a prime-field least-nonresidue bound. Its general-extension
Corollary~5.7, pp.~6--7, additionally invokes BGL Theorem~2 and AL86
for the root-of-unity construction; BGL's theorem takes certain
auxiliary irreducible polynomials as input
\cite[Theorem~2, p.~7]{CXNBGL01}.
The final broad derandomization claim in BIMS17 does not itself
supply a reduction eliminating that separate input under the narrower
premise. We therefore do not infer the unrestricted general-extension
result from $H_A$ by this citation chain. CX-N13 proves a stronger
positive transfer for every $r\mid p-1$ directly. This source-scope
conclusion makes no claim that a different $H_A$-only route is
impossible.


\section{What the supplied refutation of ETH changes}
\label{sec:eth}

\subsection{The precise quantitative premise}

OpenAI's supplied Theorem~1.1, with source label
\texttt{thm:main}, is stronger than its abstract formulation
\cite[Theorem~1.1]{OpenAIETH}.
For each fixed integer \(k\geq2\), it counts the satisfying assignments of a
\(k\)-CNF with \(n\geq20\) declared variables and \(m\) input clauses in
\[
 (n+m)^{O_k(1)}
 2^{O_k(n\log\log n/\sqrt{\log n})}
\]
bit operations and \((n+m)^{O_k(1)}\) bits of space.
For each fixed \(D>0\), if \(m\leq Dn\), the running time improves to
\[
 (n+m)^{O_{k,D}(1)}2^{O_{k,D}(n/\sqrt{\log n})}.
\]
These are theorem premises under the convention requested in this task.
No Lean replay or independent verification of the complete counting
algorithm is claimed here.

The constants, polynomial exponents, and thresholds may depend on the
\emph{fixed} width \(k\). The statement gives no bound uniform in
\(k=k(n)\). For fixed \(k\), duplicate clauses can be removed and
tautologies discarded without changing the count; the number of distinct
non-tautological clauses is at most
\(\sum_{j=0}^{k}2^j\binom nj=n^{O_k(1)}\).
Reading a repetitively encoded input still incurs its polynomial
input-length cost. In this section, \emph{subexponential} means
\(2^{o(n)}\) in the parameter explicitly named; it does not mean the
sometimes-used class
\(\bigcap_{\epsilon>0}\operatorname{DTIME}(2^{n^\epsilon})\).

\subsection{CX-E01: which exponential-time hypotheses are refuted}
\label{finding:CX-E01}

\paragraph{Finding.}
The supplied theorem refutes ETH, \(\#\mathrm{ETH}\), fixed-width SETH,
the analogous counting SETH, their bounded-error randomized versions,
Gap-ETH, and NSETH. Existing conditional lower-bound theorems remain
valid implications, but their refuted premises no longer establish
hardness. Neither side of \(\mathrm P\) versus \(\mathrm{NP}\) is
established by these deductions.

\paragraph{Proof.}
For fixed \(k\), choose a constant \(C_k\) bounding the exponential
coefficient in the theorem and put
\[
 h_k(n)=C_k n\log\log n/\sqrt{\log n}.
\]
Since \(h_k(n)/n\longrightarrow0\), the exact counter runs in
\(2^{o(n)}(n+m)^{O_k(1)}\). Its answer is positive exactly when the
formula is satisfiable and zero exactly when it is unsatisfiable.
With \(k=3\), this contradicts the positive exponential rate required
by ETH and by \(\#\mathrm{ETH}\).

SETH asserts that for every \(\epsilon>0\) some fixed \(k\) has no
algorithm of time
\(2^{(1-\epsilon)n}(n+m)^{O(1)}\).
Fix \(\epsilon=1/2\). For \emph{each} fixed \(k\),
\(h_k(n)\leq n/2\) for all sufficiently large \(n\).
Finitely many smaller values are absorbed into a \(k\)-dependent
constant. Thus every fixed width admits the required algorithm with
the same exponential coefficient \(1/2\). The negation of SETH does
not require its polynomial factor or threshold to be independent of
\(k\). The same proof applies to its counting analogue.

A deterministic algorithm is a randomized algorithm with error zero,
so the corresponding bounded-error hypotheses are also contradicted.
On the Gap-ETH promise, return YES precisely when the exact count is
positive. This distinguishes a satisfiable 3-CNF from one in which every
assignment violates a specified positive fraction of clauses. With
the usual \(m=O(n)\) promise, the fixed-density bound is directly
\(2^{o(n)}\operatorname{poly}(n)\). No approximation algorithm is
needed. Finally the zero-count test is a deterministic, hence also
nondeterministic, subexponential algorithm for fixed-width CNF
unsatisfiability. It contradicts NSETH.
The corresponding hypothesis formulations and their direct use of the
counter also appear in the supplied preprint
\cite[\S1.1.1--\S1.1.3]{OpenAIETH}.

For the logical scope, the formulas \(H\Rightarrow B\) and \(\neg H\)
entail neither \(B\) nor \(\neg B\): assigning either truth value to
\(B\) satisfies both formulas. Thus deleting the premise of a
conditional lower bound is not a construction of a faster algorithm
for its target. This observation is not an independence theorem for
any particular complexity conjecture.
The displayed running-time benchmark is superpolynomial, since
\[
 \frac{n\log\log n/\sqrt{\log n}}{\log n}\longrightarrow\infty.
\]
It is therefore not a polynomial-time bound for SAT. An upper bound
also supplies no superpolynomial lower bound for SAT. These deductions
establish neither \(\mathrm P=\mathrm{NP}\) nor
\(\mathrm P\ne\mathrm{NP}\).
\hfill\(\square\)

\paragraph{Sparsification and its overhead.}
The full fixed-width statement already applies without a sparsification
step. If a separate procedure produces at most \(2^{\epsilon n}\)
fixed-width formulas on at most \(n\) variables, each with at most
\(D_\epsilon n\) clauses, in
\(2^{\epsilon n}\operatorname{poly}_{k,\epsilon}(n+m)\) time,
pad each formula to \(n\) declared variables by forcing each new
variable to zero. This preserves its count and gives density at most
\(D_\epsilon+1\).
The fixed-density counter then gives the combined bound
\[
 \operatorname{poly}_{k,\epsilon}(n+m)
 2^{\epsilon n+O_{k,D_\epsilon}(n/\sqrt{\log n})}.
\]
This follows by multiplying the per-instance cost by the number of
instances and charging their construction.
For decision, a disjunction-equivalent list suffices.
For counting by addition, the represented satisfying-assignment sets
must partition the original set, or a separately proved counting
reconstruction must be provided: overlapping lists cannot simply
have their counts summed.
For each fixed \(\epsilon\) the displayed exponent is
\(\epsilon n+o(n)\).
Substituting \(\epsilon=\epsilon(n)\) additionally requires control of
the constants and of the reduction's uniformity.
None of these extra premises is needed for the direct use of
Theorem~1.1 above.

\paragraph{Limits of the refutation.}
The theorem does not give a
\(2^{(1-\delta)n}\operatorname{poly}(n,m)\) algorithm for
\emph{unrestricted-width} CNF-SAT, because its width may grow with \(n\).
Nor does the refutation of SETH logically refute OVC, the 3SUM or APSP
conjectures, or \(\mathrm{FPT}\ne\mathrm{W[1]}\).
An implication from SETH to a weaker conjecture has the wrong direction
for that inference. Unconditional lower bounds, NP-hardness
reductions, and results based on independent assumptions retain their
previous status.

\subsection{CX-E02: a parameter-preserving transfer criterion}
\label{finding:CX-E02}

\paragraph{Finding.}
Suppose a problem on inputs of length \(L\), with parameter \(u\),
has a polynomial-time reduction to fixed-width \(k\)-CNF with
\(N\leq au+b\) variables and \(N+M\leq L^{O(1)}\) variables and clauses
in the standard explicit encoding, where
\(a,b,k\) are fixed constants. Then it is decidable in
\[
 L^{O(1)}
 2^{O(u\log\log u/\sqrt{\log u})}
\]
time and polynomial space.
A parsimonious reduction gives the same bound for its witness count.
If the formula additionally has \(M=O(N)\) clauses, the
\(\log\log u\) factor is absent. Polynomially many such calls and
polynomial-time answer processing preserve these bounds.

\paragraph{Proof.}
Construct the formula and run the supplied counter or its positivity
test. For \(N\leq au+b\), after changing constants to cover finitely
many small values,
\[
 N\log\log N/\sqrt{\log N}
 =O(u\log\log u/\sqrt{\log u}).
\]
A polynomial-time reduction produces at most polynomially many bits,
so the prefactor is \(L^{O(1)}\).
Parsimoniousness identifies the two counts exactly.
The fixed-density part of Theorem~1.1 gives the stated improvement
when \(M=O(N)\).
Polynomially many calls change only the polynomial factor; their
workspaces can be reused.
\hfill\(\square\)

A reduction with \(N=\Theta(u^c)\), \(c>1\), instead gives a benchmark
exponent of order \(u^c\log\log u/\sqrt{\log u}\), which is not \(o(u)\).
Ordinary polynomial-time completeness therefore does not suffice to
transfer subexponential time in the original parameter.
Similarly, a \(2^{o(v)}\) graph algorithm is not a bound
\(f(k)v^{O(1)}\) in solution size \(k\), and a formula with \(O(v)\)
variables does not provide a \(2^{o(t)}\operatorname{poly}(v)\)
algorithm in treewidth \(t\).

For a claimed polynomial-exponent improvement, such as a
subquadratic algorithm in \(L\), the actual reduction cost,
solver prefactor, number of calls, and answer-processing cost must
each be bounded within that target time.
The notation \(L^{O(1)}\) alone cannot certify this.
For fixed-parameter tractability in a parameter \(t\), a sufficient
separate reduction would have fixed clause width,
\(N\leq g(t)\), and \(M\leq h(t)L^C\), and run in
\(h(t)L^C\) time, for computable \(g,h\) and a constant \(C\).
Substitution into Theorem~1.1 then gives
\(F(t)L^{O(1)}\) for a computable \(F\): every factor involving
the bounded variable count is absorbed into \(F(t)\).
The theorem does not itself construct such a reduction.

The supplied Corollary~1.3 already gives exact subexponential counting
of independent sets and fixed-\(q\) graph colorings in the vertex
parameter, and of perfect matchings in \(v+e\)
\cite[Corollary~1.3]{OpenAIETH}.
Bounded-set-size Set Splitting has a direct encoding with one bit per
universe element and, for each set, an all-positive and an all-negative
clause. These are actual algorithmic transfers.
Deleting a SETH premise alone supplies no subquadratic algorithm for
Edit Distance, LCS, or arbitrary Orthogonal Vectors.

\paragraph{A further explicit parameter check: counting cliques.}
For a graph with \(v\) vertices, the number of \(k\)-cliques can be
computed in
\(\operatorname{poly}(v+e)
 2^{O(v\log\log v/\sqrt{\log v})}\),
uniformly for \(0\leq k\leq v\).
Use a selection bit for every vertex, and for each nonedge
\(\{a,b\}\) impose \(\neg x_a\vee\neg x_b\).
Add a Boolean circuit requiring exactly \(k\) selected vertices.
A balanced tree of binary additions computes the population count
with \(O(v)\) gates: level \(j\) has \(O(v/2^j)\) additions on
\(O(j)\) bits, and
\(\sum_{j\geq1}vj/2^j=O(v)\).
Comparison with \(k\) adds \(O(\log v)\) gates.
The gate encoding in CX-E03 gives \(O(v)\) variables and \(O(v^2)\)
clauses, with a unique extension for each clique.
Theorem~1.1 proves the bound. This does not assert fixed-parameter
tractability in \(k\).
\hfill\(\square\)

\subsection{CX-E03: the Circuit-SAT saving and its exact parameter}
\label{finding:CX-E03}

\paragraph{Finding.}
An arbitrary Boolean circuit with \(n\) inputs, \(s\) noninput gates
of fan-in at most two, and unrestricted fan-out has a deterministic
satisfiability and exact accepting-input-count algorithm of time
\[
 (n+s)^{O(1)}
 2^{O((n+s)/\sqrt{\log(n+s)})}
\]
and polynomial space.
For \(s=O(n)\) this is \(2^{o(n)}\).
This bound does not meet the all-polynomial-size Circuit-SAT premise
of Williams's Theorem~1.1 \cite[Theorem~1.1]{Williams2013}.

\paragraph{Proof.}
Introduce one variable for each input and each gate.
A two-input gate has eight assignments to its inputs and output,
exactly four inconsistent with its truth function.
Forbid each inconsistent assignment by a clause of width at most
three; constants and unary gates need no more clauses.
Add a unit clause setting the final output to one.
The resulting formula has \(N=n+s\) variables and
\(M\leq4s+1\leq5N\) clauses.
Every input assignment extends to exactly one assignment to the gate
variables, by evaluation in topological order.
Thus this reduction preserves the number of accepting inputs.
Apply the fixed-density part of Theorem~1.1 with \(k=3,D=5\).
\hfill\(\square\)

Precisely, Williams's Theorem~1.1 assumes a superpolynomial function
\(q(n)\) such that Circuit-SAT on \(n\) inputs and \(n^a\) gates
has a deterministic or co-nondeterministic algorithm of time
\[
 \frac{2^n\operatorname{poly}(n^a)}{q(n)}
 \qquad\text{for every fixed }a.
\]
Its conclusion is \(\mathrm{NEXP}\not\subseteq\mathrm{P/poly}\).
For \(a>1\), substituting \(s=n^a\) into the proved benchmark gives an
exponent of order \(n^a/\sqrt{\log n}>n\).
Consequently this reduction does not supply that premise.
This is a diagnosis of the particular attempted inference, not an
impossibility theorem for other approaches to stronger lower bounds.

\subsection{CX-E04: the quantitative Boolean circuit lower bound}
\label{finding:CX-E04}

\paragraph{Finding.}
Under the task's theorem convention there exist a language
\(L\in\mathrm E^{\mathrm{NP}}\) and \(c>0\) such that
\[
 C_L(n)>c\,n\sqrt{\log n}
 \quad\text{for infinitely many }n,
\]
where \(C_L\) counts arbitrary fan-in-two Boolean gates.
In particular,
\(\mathrm E^{\mathrm{NP}}\not\subseteq\mathrm{SIZE}(O(n))\).
The quantitative statement is already
\cite[Corollary~1.4]{OpenAIETH}; it is not attributed as a new
circuit lower-bound discovery to this run.
The qualitative implication is
\cite[Theorem~6.1]{Williams2013}.
The finding here is the exact transfer and its scope.

\paragraph{Proof of the transfer.}
The decision algorithm required by the supplied Corollary~1.4 is
obtained from the fixed-density part of Theorem~1.1 by testing the
count for positivity.
Its premise is therefore satisfied, independently of both GRH
sources. Apply that corollary.
Since \(\sqrt{\log n}\longrightarrow\infty\), its language cannot
have a circuit family of eventual size \(O(n)\).
\hfill\(\square\)

\paragraph{Check of the cited implication proof.}
The complete proof in the supplied
\nolinkurl{quantitative-circuits.tex}, inserted in Appendix~C.1, was
read and its quantitative parameters checked.
It uses the unary nondeterministic time hierarchy in
\cite[Theorem~2.1]{Williams2013} and the local reduction
\cite[Theorem~1]{JMV2013}.
The latter gives a clause decoder with
\(q=n+O(\log n)\) index bits and \(O(n)\) output bits, each depending
on only constantly many index bits.
Enumerating inputs of bounded Hamming weight recovers supports and
truth tables, hence an \(O(n)\)-gate decoder in polynomial time.
The bits of the first accepting assignment, indexed by padded
\(dn\)-bit variable names, form a language in
\(\mathrm E^{\mathrm{NP}}\).
A hypothetical witness circuit of size
\[
 s_n=\left\lfloor\eta n
          \left\lfloor\sqrt{\log_2 n}\right\rfloor\right\rfloor
\]
is guessed. Three copies together with the decoder test whether a
clause is violated.
The gate encoding produces
\[
 U_n=(3\eta+o(1))n\sqrt{\log_2 n}
\]
variables and at most \(5U_n\) clauses.
For a solver bound \(V^b2^{AV/\sqrt{\log_2 V}}\), its cost is
\(2^{(3A\eta+o(1))n}\).
Choosing \(\eta<1/(12A)\) leaves enough margin, including polynomial
overhead, to contradict the chosen
\(\mathrm{NTIME}(2^n)\setminus\mathrm{NTIME}(2^{n/2})\)
language.
No gap was found in this implication proof.
This is not a claimed reproof of the counter, the hierarchy theorem,
or the local-reduction theorem.

The result is a bound in \(\mathrm E^{\mathrm{NP}}\), not in NP;
it concerns nearly linear Boolean circuits, not all polynomial-size
circuits or arithmetic circuits for the permanent.
The displayed existential constant \(c\) and infinitely-often
quantifier should not be strengthened to
\(\mathrm E^{\mathrm{NP}}\not\subseteq
  \mathrm{SIZE}(O(n\sqrt{\log n}))\).

\subsection{CX-E05: a quantitative sparse permanent algorithm}
\label{finding:CX-E05}

\paragraph{Finding.}
For \(A\in\{0,1\}^{r\times r}\) with \(e\) nonzero entries, its
integer permanent can be computed in
\[
 \operatorname{poly}(r+e)
 2^{O((r+e)/\sqrt{\log(r+e)})}
\]
time and polynomial space.
When \(e=O(r)\), the bound is
\(\operatorname{poly}(r)2^{O(r/\sqrt{\log r})}\).
This makes quantitative the matching consequence of the supplied
Corollary~1.3 by using its sharper fixed-density Theorem~1.1.

\paragraph{Proof.}
Interpret \(A\) as a bipartite graph with \(2r\) vertices and \(e\)
edges, with one selection bit for each edge.
At a vertex of degree \(d\), maintain bits \(a_i,b_i\) saying that
among the first \(i\) incident edges, respectively, at least one and
at least two have been selected. Initialize \(a_0=b_0=0\), and set
\[
 a_i=a_{i-1}\vee x_i,\qquad
 b_i=b_{i-1}\vee(a_{i-1}\wedge x_i).
\]
Induction on \(i\) proves these meanings.
Exactly one incident edge is selected if and only if
\(a_d\wedge\neg b_d=1\).
This requires \(O(d+1)\) gates.
Summing degrees, and conjoining the tests, gives \(O(r+e)\) gates;
an isolated vertex may instead be rejected immediately.

Apply the parsimonious encoding of CX-E03.
It has \(N=O(r+e)\) variables and \(M=O(N)\) clauses.
Every perfect matching gives one satisfying assignment with uniquely
determined gate bits, and every satisfying assignment selects a
perfect matching.
The number of these matchings is
\[
 \operatorname{per}(A)=
 \sum_{\sigma\in S_r}\prod_{i=1}^r A_{i,\sigma(i)}.
\]
The fixed-density bound proves the assertion.
\hfill\(\square\)

For dense matrices this encoding has \(\Theta(r^2)\) variables; its
bound does not give \(2^{o(r)}\) time in matrix dimension.
Nor does it construct an arithmetic circuit for the formal
polynomial \(\operatorname{per}_r(X)\).
The algorithm evaluates an integer count at a supplied 0--1 matrix.

\subsection{CX-E06: VP, VNP, and the surviving GCT gates}
\label{finding:CX-E06}

\paragraph{Finding and exact logical scope.}
Neither \(\mathrm{VP}=\mathrm{VNP}\) nor
\(\mathrm{VP}\ne\mathrm{VNP}\), and neither a polynomial arithmetic
circuit nor a superpolynomial arithmetic lower bound for the
permanent, is established by CX-E01--CX-E05.
The deductions do not discharge the border-hardness or black-box
identity-testing assumptions in GCT V Theorems~1.7, 2.4, 5.11,
and~5.14. This is a statement about the premises of these particular
implications, not a formal independence theorem.

\paragraph{Dictionary and proof.}
Over a fixed field, VP consists of polynomial-degree families with
polynomial-size arithmetic circuits. VNP permits Boolean-cube sums
of such polynomial families.
For characteristic other than two, the permanent is VNP-complete
under polynomial projections
\cite[Theorem~2, p.254]{Valiant1979}.
The determinant is complete for VBP, the weakly skew or branching
program subclass; it is not thereby complete for all of VP.
The corresponding orbit-closure formulation uses the border of
that model
\cite[Definitions~2.9, 2.25; Theorems~2.20, 2.29, 7.3]{Burgisser2024}.

The supplied counter accepts a finite bit-string description of a
fixed-width CNF and returns an integer.
Its conclusion does not construct, from an arithmetic verifier
\(g_n(X,Y)\), an arithmetic circuit for the symbolic polynomial
\[
 \sum_{Y\in\{0,1\}^{p(n)}}g_n(X,Y).
\]
CX-E05 also evaluates only a numerical specialization.
The lower bound in CX-E04 has a language in
\(\mathrm E^{\mathrm{NP}}\) as its target and Boolean gates as its
model. It has neither the family nor the model required for an
arithmetic permanent lower bound.
Thus these deductions do not furnish the premises of the cited
algebraic implications.

There is a separate quantitative obstruction to a tempting inference.
GCT V Theorem~1.7, source label \texttt{tintronnldetnew}, assumes
that for some fixed \(\epsilon>0\) the permanent cannot be approximated
infinitesimally by affine-linear symbolic determinants of size
at most \(2^{r^\epsilon}\).
Even if a new argument provided an upper bound
\(2^{O(r/\sqrt{\log r})}\) in that same determinant or border model,
it would not refute
\(2^{r^\epsilon}\) hardness for \(0<\epsilon<1\), because
\[
 \frac{r/\sqrt{\log r}}{r^\epsilon}
 =\frac{r^{1-\epsilon}}{\sqrt{\log r}}\longrightarrow\infty.
\]
No such symbolic conversion has been established here in the first
place.
The remaining gate is a substantive arithmetic conversion or
obstruction construction, not an erased GRH assumption.
\hfill\(\square\)

The GCT assumptions just mentioned are exactly located in the supplied
source as \texttt{tintronnldetnew}, \texttt{trussell2},
\texttt{tnnlexpfull}, and \texttt{tequisdit}, respectively.
Improving the complexity of finding a normalization system is also
distinct from proving the orbit-closure noncontainment required by a
permanent lower bound.



\clearpage\part{Checked geometric restatements and dictionaries}

% CX_GRH: Part II geometry fragment, separate from the complexity section.
% Root supplies the exact published/archival locator crosswalk.

\section{CX-K08: arithmetic fibers and the geometry of normalization}
\label{cx:k08}

\paragraph{Statement.}
The Koiran reduction describes arithmetic fibers of an affine
finite-type $\mathbb Z$-scheme and a finite \'etale permutation local
system on an open subscheme of $\operatorname{Spec}\mathbb Z$.
The h.s.o.p. search in CX-K05 finds a finite morphism to affine space;
for positive-dimensional standard graded rings, its projective form
is a finite basepoint-free linear projection. These dictionaries
preserve the proofs and their analytic hypotheses.

\begin{proof}
Given integer equations $f_i$, put
\[
 X=\operatorname{Spec}\bigl(\mathbb Z[x_1,\ldots,x_t]/(f_i)\bigr).
\]
A complex solution exists exactly when $X_{\mathbb Q}$ is nonempty.
Indeed, a complex point makes the rational ideal proper; a maximal
ideal of the resulting finitely generated $\mathbb Q$-algebra then
has a finite extension $E/\mathbb Q$ as residue field. Conversely,
such a finite residue field embeds in $\mathbb C$. The closed point
defines $\operatorname{Spec}E\to X_{\mathbb Q}$. Clearing the finitely
many coordinate denominators extends it to
\[
 \operatorname{Spec}\mathcal O_E[1/B]
     \longrightarrow X_{\mathbb Z[1/B]}
\]
for some nonzero integer $B$. Every residue-degree-one closed point
above a rational prime $p$ gives an $\mathbb F_p$-point of $X$.
These are the specialization maps counted in CX-K01.

Choose the integral primitive element $\beta$ used there and enlarge
$B$ to include the index $[\mathcal O_E:\mathbb Z[\beta]]$.
For a fixed auxiliary prime $\ell$, set
\[
 U=\operatorname{Spec}\mathbb Z[1/(\ell B D_E)].
\]
The normalization $\pi:\widetilde U\to U$ in $E$ is finite
\'etale. Its direct image $\pi_*\mathbb Q_\ell$ is a rank
$[E:\mathbb Q]$ permutation local system. At each $p\in U$,
Frobenius acts on the embeddings of the residue factors by
permutation. Its trace counts fixed points, hence residue-degree-one
primes above $p$. Because the index is inverted, this also equals the
number of roots of the chosen integral primitive-element polynomial
in $\mathbb F_p$.

The Euler product of this local system is $\zeta_E$ with the finitely
many excluded Euler factors removed. The Frobenius eigenvalues are
roots of unity, so the local system is a weight-zero Artin object.
That local purity fact does not determine the global zero set of its
number-field Euler product. In particular, it does not eliminate the
analytic premise in CX-K01.

If $E/\mathbb Q$ is abelian, its permutation representation decomposes
into rank-one characters after extending coefficient scalars to a
splitting field. This gives the Dirichlet factorization used in
CX-K03. For arbitrary $E$, the representation is
\[
 \operatorname{Ind}_{G_E}^{G_{\mathbb Q}}\mathbf1,
\]
which need not decompose into one-dimensional representations.
In a representation-ring dictionary it is the class of this
permutation object in $K_0$. Applying $\lambda$-operations to that
class does not supply a theorem about global zeros of its $L$-function.
This identifies the representation-theoretic object without adding
an Artin-holomorphy assumption to the proof.

Now let $R$ be as in CX-K05 and let $d=\dim R$. The condition that
$L_1,\ldots,L_d$ vanish only at the vertex is equivalent, by the
finite-module argument there, to finiteness of
\[
 \operatorname{Spec}R\longrightarrow\mathbb A_{\mathbb Q}^{d}.
\]
For $d\ge1$, the forms have no common zero on $\operatorname{Proj}R$.
They define the finite linear projection
\[
 \operatorname{Proj}R\longrightarrow\mathbb P_{\mathbb Q}^{d-1}.
\]
For $d=0$, $\operatorname{Proj}R$ is empty, so no such projective
target is asserted. Finiteness in positive dimension follows from the
graded finite-module presentation, or by localizing it on the
standard target charts and taking degree-zero pieces.

Candidate matrices of coefficients form an affine parameter scheme
$T$. The incidence condition that a projective point lie in all
candidate kernels is closed in
$\operatorname{Proj}R\times_{\mathbb Q}T$.
The projection to $T$ is proper, since $\operatorname{Proj}R$ is
projective over $\mathbb Q$, so its image is closed. Its complement
is exactly the open locus of successful matrices. The bounded
coefficient argument and oracle search in CX-K05 are effective
statements about rational points of that open locus, with the stated
bit bounds.

The finite morphism also spreads out. Finiteness over $\mathbb Q$
provides finitely many monic integral-dependence relations for the
original generators $z_i$, with coefficients in $\mathbb Q[L]$.
Clear denominators in the $L_i$, in these coefficients, and in their
finite expressions modulo the defining ideal. After inverting a
single nonzero integer $B$, the relations hold over
$\mathbb Z[1/B]$ and show that the corresponding morphism is finite
over that base.

This last assertion is precisely finiteness after localization. It
does not assert flatness at every remaining prime or that every
special fiber is a Noether normalization with an algebraically
independent parameter system of the same dimension. No
Hilbert- or Quot-scheme representability theorem is strengthened by
the complexity bound. The mathematical dictionary therefore preserves
both the analytic and geometric limits of the preceding deductions.
\end{proof}

The archival and published EGA/SGA/FGA locator crosswalk is supplied
separately for the constructions actually used: finite morphisms,
integrality, spreading out, projective properness, and finite \'etale
permutation objects. No additional conclusion about motives,
$\gamma$-filtrations, Schur functors, or global zeros follows merely
from placing the same objects in those languages.

This full proof was reread by \texttt{/root/koiran\_gct}, GPT-6 Astra
Pro, on 20 November 2026 UTC, with inherited effort setting not
exposed. It was also read in full by \texttt{/root/eth\_consequences}
and \texttt{/root}. The index localization, exclusion of $\ell$,
and $d=0$ projective case were explicitly checked and incorporated.


\section{The ETH findings as scheme and Grothendieck-group statements}
\label{sec:eth-geometric}

\subsection{Geometric restatements of CX-E01--CX-E06}

The constructions and proofs in CX-E07--CX-E08 supply the following
dictionary. The input encoding, variable-count parameter, and
algorithmic bounds remain those in Part~I.

\begin{longtable}{@{}p{0.12\linewidth}p{0.83\linewidth}@{}}
\hline
Finding & Exact geometric restatement\\
\hline
CX-E01 &
For the succinct Boolean scheme \(X_F\), satisfiability is
nonemptiness and the satisfying-assignment count is its finite
étale degree. The theorem computes that degree in the stated
subexponential time in the number of Boolean coordinates.\\[4pt]
CX-E02 &
A decision reduction preserves emptiness; a parsimonious reduction
preserves the rank of the corresponding constant finite étale
scheme. The quantitative transfer requires the stated linear bound
on Boolean coordinates. A witness bijection also induces an
isomorphism of these constant schemes.\\[4pt]
CX-E03 &
The gate scheme projects isomorphically to the scheme of accepting
input assignments: every auxiliary coordinate is a uniquely
determined Boolean polynomial in the earlier coordinates.
Thus adding gate variables preserves both nonemptiness and degree.\\[4pt]
CX-E04 &
The degree algorithm yields the cited lower bound for Boolean
circuits of a language in \(\mathrm E^{\mathrm{NP}}\).
Its circuit size counts Boolean gates; the dictionary does not
identify these circuits with a representation of an algebraic
coordinate ring or with arithmetic circuits.\\[4pt]
CX-E05 &
The edge-selection constraints define a constant finite étale
perfect-matching scheme. Its degree is the permanent of the
specified \(0\)--\(1\) matrix; the construction has \(O(r+e)\)
Boolean coordinates and a fixed-density presentation.\\[4pt]
CX-E06 &
The computed integer is a rank of a finite free algebra.
A symbolic permanent family and a graded equivariant
orbit-closure coordinate ring require further data.
The rank statement therefore supplies neither the arithmetic
circuit construction nor the Schur-multiplicity obstruction
specified by those algebraic questions.\\
\hline
\end{longtable}

\subsection{CX-E07: counting as the degree of a finite étale scheme}
\label{finding:CX-E07}

\paragraph{Finding.}
For a Boolean \(k\)-CNF \(F\), put
\[
 X_F=\operatorname{Spec}
 \frac{\mathbb Z[x_1,\ldots,x_n]}
 {\bigl(x_i^2-x_i\ (1\leq i\leq n),\;
     \prod_{\ell\in C}(1-\ell(x))\ (C\text{ a clause of }F)\bigr)},
\]
where \(\ell(x)=x_i\) for a positive literal and
\(\ell(x)=1-x_i\) for a negative literal.
There is an isomorphism
\[
 X_F\cong\bigsqcup_{\#F}\operatorname{Spec}\mathbb Z.
\]
Consequently \(p:X_F\to\operatorname{Spec}\mathbb Z\) is finite étale,
and
\[
 \operatorname{rank}_{\mathbb Z}p_*\mathcal O_{X_F}=\#F,
 \qquad
 [p_*\mathcal O_{X_F}]=\#F[\mathcal O]\quad\hbox{in }K_0(\mathbb Z).
\]
The supplied theorem computes this degree, from the displayed
bounded-locality presentation, with its stated time and space bounds.
The presentation is succinct; an explicit list of the components or
a basis of its coordinate algebra need not be produced.

\paragraph{Proof.}
The ideals \((x_i)\) and \((x_i-1)\) are comaximal in
\(\mathbb Z[x_i]\).
The Chinese remainder theorem, applied successively in all variables,
identifies
\[
 \mathbb Z[x_1,\ldots,x_n]/(x_i^2-x_i)_i
 \cong\prod_{a\in\{0,1\}^n}\mathbb Z
\]
by evaluation.
The polynomial attached to a clause evaluates to one at an assignment
violating that clause, and to zero at one satisfying it.
Quotienting by all clause polynomials therefore removes precisely
the factors corresponding to nonsatisfying assignments.
The surviving factors give the asserted disjoint-union isomorphism.
Their coordinate module is free of rank \(\#F\), which proves the
rank and \(K_0\) identities.
The input formula and the factored clause presentation use the same
literal data, so the counter computes that rank without changing the
parameters beyond fixed encoding constants.
\hfill\(\square\)

\paragraph{The exact dictionary for CX-E01--CX-E05.}
Satisfiability is nonemptiness of \(X_F\).
A witness chooses one of its components.
Counting computes the degree of \(p\).
A parsimonious reduction in CX-E02 preserves the number of components
and hence the isomorphism class of the constant finite étale scheme.
A specified witness bijection identifies the component sets.
In the gate constructions of CX-E03 and CX-E05 there is a more
explicit isomorphism: projection forgets gate variables, and its
inverse assigns each gate the multilinear Boolean polynomial of
its predecessors.
For each Boolean input, the gate clauses permit exactly that value;
the product decomposition proved above therefore shows that the
projection and its inverse compose to the identity as morphisms
over \(\mathbb Z\).
The gate variables introduce no extra components.
All quantitative bounds remain bounds in the size of the specified
presentation; they are not complexity assertions about a scheme
given only up to an unencoded isomorphism.

\paragraph{Exterior powers, weights, and the exact limits.}
In the \(\lambda\)-ring of vector bundles,
\[
 \lambda_t([p_*\mathcal O_{X_F}])=(1+t)^{\#F},
\]
because the \(j\)-th exterior power of a free module of rank \(\#F\)
has rank \(\binom{\#F}{j}\).
Its reduced rank class is zero, so this construction produces no
nonzero class in the positive \(\gamma\)-filtration.
For a fiber over \(\mathbb F_q\) and \(\ell\ne\operatorname{char}\mathbb F_q\),
the geometric \(\ell\)-adic cohomology is concentrated in degree zero,
has dimension \(\#F\), and has identity Frobenius.
The fiber's zeta function is consequently
\((1-t)^{-\#F}\).
These facts follow component by component from the proved
decomposition. The object has weight zero and its Artin motive is
a sum of unit objects; no new purity theorem is asserted.

\subsection{CX-E08: the additional data a GCT obstruction would require}
\label{finding:CX-E08}

\paragraph{Finding.}
The rank computation in CX-E07 does not, through this dictionary,
construct a separating representation in a determinant-orbit-closure
coordinate ring.

\paragraph{Exact comparison and proof.}
Let \(E=\mathbb Z^{m^2}\), let \(G=\mathrm{GL}(E)\), and let
\(\mathbb V\) be the affine coefficient space associated to
\(\operatorname{Sym}^m(E^*)\).
Take the schematic closure in \(\mathbb V\) of the determinant orbit.
After base change to the complex numbers, a containment of a padded
permanent orbit closure \(Y\) in the determinant orbit closure \(X\)
would give a graded \(G\)-equivariant quotient
\[
 \mathbb C[X]\longrightarrow\mathbb C[Y].
\]
In every finite homogeneous degree, rational representations of this
general linear group are semisimple.
Write \(E_{\mathbb C}=E\otimes_{\mathbb Z}\mathbb C\).
For each irreducible polynomial representation
\(S_\lambda E_{\mathbb C}\), equivariance
therefore sends its isotypical source summand onto the target
summand. Writing the two summands as
\(S_\lambda E_{\mathbb C}\otimes M_\lambda\) and
\(S_\lambda E_{\mathbb C}\otimes N_\lambda\), surjectivity implies
\(\dim N_\lambda\leq\dim M_\lambda\).
A reversed inequality would obstruct the containment
\cite[\S7.4, (7.6), p.~38, and (7.9), p.~39]{Burgisser2024}.
This is a statement in the Grothendieck group of rational
representations, with Schur and exterior-power operations, rather
than a statement about total rank alone.

CX-E07 provides a bounded-degree Boolean presentation and the integer
\(\#F\).
Its decomposition into unit components supplies no such graded
\(G\)-representation, map of coordinate rings, or inequality of Schur
multiplicities.
Consequently this proved dictionary does not produce that obstruction.
A new construction furnishing these additional data, with the
appropriate complexity guarantees, would be needed.
This verifies the missing data in the attempted implication; it does
not rule out a future construction.
\hfill\(\square\)

The exact EGA, SGA, FGA and archival locators for the constructions in
this section are given in the paper's common geometric source
dictionary. The mathematical identifications above use only the
displayed decomposition, finite-module operations, and the
characteristic-zero representation statement explicitly indicated.


\section{A common geometric dictionary, with integral models}
\label{sec:common-geometry}

Every restatement in this part retains the input encoding and the hypotheses
of Part~I.  A complexity theorem is not an isomorphism-invariant assertion
about an unencoded abstract scheme.  Published EGA/SGA/FGA locators support
the mathematical constructions.  The archive cotes identify inspected
historical companions; they are not presented as manuscripts containing AM,
ETH, or GCT complexity assertions.

\subsection{CX-G01: the representation-ring form of the analytic gates}
\label{finding:CX-G01}

\textbf{Finding.} Let $U=\operatorname{Spec}\mathbb Z[1/(B\ell)]$ and let
$\pi:Y\to U$ be the finite étale cover obtained by spreading out a number
field $E/\mathbb Q$, with all ramified primes inverted.  Its permutation
local system has class
\[
 [\pi_*\mathbb Q_\ell]
 = [\operatorname{Ind}_{G_E}^{G_{\mathbb Q}}\mathbf1]
 \quad\text{in }K_0(\operatorname{Rep}_{\mathbb Q_\ell}^{\mathrm{fin}}G_{\mathbb Q}).
\]
Its unramified Euler product is $\zeta_E$ with the excluded local factors
removed.  This exactly transports CX-R01--R02 and CX-K01--K04 into the
language of finite étale covers and representation Grothendieck groups.

\begin{proof}
The geometric fiber is the finite continuous $G_{\mathbb Q}$-set
$\operatorname{Hom}_{\mathbb Q}(E,\overline{\mathbb Q})$.
The stalk of the pushforward is the vector space of functions on that set.
A Frobenius cycle of length $f$ has characteristic polynomial $T^f-1$,
and hence reciprocal local factor $1-t^f$.  Its cycle lengths are the
residue degrees of the primes of $E$ above $p$.  Taking $t=p^{-s}$ gives
exactly the factors $(1-p^{-fs})$ in the reciprocal of the Dedekind Euler
factor.  Trace counts fixed points, hence degree-one primes, as needed by
Koiran.  The equivalence between finite étale covers and finite
fundamental-group sets is \cite[Exp.~V, \S4 Theorem~1, and \S7]{SGA1}.

For an abelian $E/\mathbb Q$, character decomposition gives the primitive
Dirichlet factorization in CX-R01.  For the cubic field of CX-R02, the
commutator trace calculation proves that the class is not a virtual sum of
rank-one objects, even allowing other conductors.  Taking $K_0$ does not
remove this obstruction.  The nonabelian Dedekind hypothesis remains the
same conditional input $M$.
\end{proof}

A finite-order ray-class character over $K$ instead defines a rank-one
local system on an open of $\operatorname{Spec}\mathcal O_K$, through
class field theory \cite[Introduction, Theorems~0.4--0.5]{MilneCFT}.
Its pushforward to an open of $\operatorname{Spec}\mathbb Z$ generally has
higher rank.  This is the exact dictionary for the Hecke hypotheses in
CX-N05--N08 and the class-field applications.  The corresponding objects
are Artin motives, with a coefficient field containing the finite character
values.  No general motivic conjecture is used.

\paragraph{Text and archive locations.}
Equivariant modules are treated in \cite[Exp.~I 4.7.1--4.7.2]{SGA3}.
The local determinant formalism is
\cite[Exp.~XV, \S3 no.~1, Proposition~1, and no.~2]{SGA5}; that source's
global rationality statement is over finite fields, not over
$\operatorname{Spec}\mathbb Z$.
The archival companions are cote~98, scan~12 (numbered sheet~11),
``Exemple I'' on finite $G$-sets, and scan~11 (sheet~10), the degree axioms
for norm morphisms \cite{Archive98}.

\subsection{CX-G02: multiplicative groups, Kummer torsors, and factoring}
\label{finding:CX-G02}

\textbf{Finding.} The modular operations of CX-N01--N06 are operations
on $\mathbb G_m$ and on explicit finite étale algebras.
An $r$-th power nonresidue $a\in\mathbb F_q^\times$, with $r\mid q-1$,
is precisely a nonzero class in
$\mathbb F_q^\times/(\mathbb F_q^\times)^r$, equivalently the Kummer
torsor
\[
 \operatorname{Spec}\mathbb F_q[t]/(t^r-a)
\]
has no $\mathbb F_q$-section.  A primitive root generates
$\mathbb G_m(\mathbb F_q)$; a generating set is a different output.

\begin{proof}
$\mathbb G_m=\operatorname{Spec}\mathbb Z[t,t^{-1}]$ represents units,
so its multiplication and power maps are literally the modular operations.
Miller's gcd tests take place in the same explicit residue ring
$\mathbb Z/N\mathbb Z$.
For the torsor, $r$ is invertible and $a\ne0$, so the derivative
$rt^{r-1}$ is a unit in the quotient.  The algebra is finite étale and
$\mu_r$ acts by multiplying $t$.  After adjoining a root, all other roots
differ from it by $\mu_r$, proving the torsor property.  A rational section
is exactly an $r$-th root of $a$.  Cyclicity gives the original power test
$a^{(q-1)/r}\ne1$ for prime $r$.  Primary-component assembly and the order
tests in Part~I are unchanged; the dictionary supplies no factorization of
the integer $q-1$.

For squarefree $f\in\mathbb F_q[t]$, the algebra
$A=\mathbb F_q[t]/(f)$ is finite étale.  Its product decomposition into
fields is the connected-component decomposition of $\operatorname{Spec}A$.
The Chinese remainder theorem sends a monic irreducible factor to its
primitive projector $e$.  Conversely, $(1-e)A$ has quotient the corresponding
field, and its inverse image in $\mathbb F_q[t]$ is generated by that factor.
Both conversions use polynomial arithmetic and linear algebra.  Repeated
factors are treated by squarefree decomposition, including Frobenius-root
extraction over the perfect base field.  Thus factoring is the effective
idempotent decomposition problem.  Existence of the finite étale
decomposition does not compute its idempotents or change the bound of the
selected factoring algorithm.
\end{proof}

\paragraph{Locations.}
$\mathbb G_m$ and $\mu_r$ occur in the diagonalizable-group formalism
\cite[Exp.~VIII 1.1--1.2]{SGA3}; the exact torsor sequence is
\cite[Exp.~IX 3.2]{SGA4}.  See also \cite[Exp.~V \S7]{SGA1}.
Cote~28, scan~3 (sheet~2), writes the action map $G\times X\to X$ and its
axioms \cite{Archive28}; cote~98, scan~12 (sheet~11), supplies the finite
$G$-set companion \cite{Archive98}.  These are locations of the language,
not historic primality or factoring algorithms.

\paragraph{The extension refinements in this language.}
CX-N13 constructs a nontrivial class in
$H^1_{\mathrm{et}}(\operatorname{Spec}\mathbb F_{p^m},\mu_r)$
for a specified prime $r\mid p-1$ from the supplied field model.
Its odd-$r$ Frobenius eigenspace is computed on that finite free
$\mathbb F_p$-algebra; its power test is exactly the torsor's
failure to have a section. The even-degree quadratic case retains
its unconditional construction.
For CX-N14, a supplied automorphism of a constant finite étale
$\mathbb F_p$-scheme gives the same component action. The moved-component
ideal and the descended idempotent construct a proper open-and-closed
subscheme. The cases $p\le d$, $r\mid p-1$, and
$v_r(p^{\operatorname{ord}_r(p)}-1)=1$, and the optional coefficient
Teichmüller unit, retain precisely their hypotheses and bounds.
For CX-N15, the Frobenius-fixed algebra
$B_f=\ker(a\mapsto a^p-a)$ inside
$A=\mathbb F_q[X]/(f)$ is $\mathbb F_p^t$, where $t$ is the number
of connected components of the finite étale scheme $\operatorname{Spec}A$.
Its primitive idempotents are exactly those of $A$.
The coefficient extension, eigenspace and equivariant interpolation
construct the unit needed to split this constant component scheme.
The resulting open-and-closed decompositions of $\operatorname{Spec}A$
therefore compute the factors in
$(d^{\log d}\log q)^{O(1)}$ time, under the exact cyclotomic--Kummer
Dedekind-zeta family in CX-N15. The pointwise Frobenius calculation is
an algebraic construction; its required number-field prime bound is
still the stated conditional analytic input.
Thus the Kummer and descent operations carry the complete proofs,
not merely their existence conclusions, into this dictionary.

\subsection{CX-G03: ideal-class computation as Picard-group computation}
\label{finding:CX-G03}

\textbf{Finding.} For $X=\operatorname{Spec}\mathcal O_K$,
the class group is $\operatorname{Pic}(X)$.  CX-N07 bounds the prime-divisor
line bundles generating this group, with exactly the same conditional
premise.  CX-N08 and the concrete class-group computation concern completeness
of the principal-divisor and unit data, with their original search bounds.

\begin{proof}
A nonzero fractional ideal $I$ of the Dedekind domain $\mathcal O_K$ is an
invertible module.  Tensor product corresponds to ideal multiplication.
Conversely, choose a basis of the generic fiber of a rank-one locally free
module and embed it in $K$; clearing denominators identifies it with a
fractional ideal.  An isomorphism of two such modules extends generically
to multiplication by a scalar in $K^\times$.  Thus isomorphism classes
are exactly fractional ideals modulo principal fractional ideals.
The prime ideal $\mathfrak p$ corresponds, in effective-divisor notation,
to $\mathcal O_X(-\mathfrak p)$; changing all signs does not affect generation.
The factor-base map is the divisor-to-line-bundle map, with kernel the
principal-divisor relations supported on that base.  Units are the kernel
of the principal-divisor map from $K^\times$.  These are precisely the
objects in the exact sequence behind CX-N08, so its index identity and
its separate relation-search cost are preserved.
\end{proof}

The Hilbert-class-field cover is
$\operatorname{Spec}\mathcal O_{H_K}\to\operatorname{Spec}\mathcal O_K$,
finite étale at finite places.  For ray class fields, delete the modulus
primes and retain the real-place conditions in the arithmetic character data.
Their characters give the local systems of CX-G01.  The required analytic
factor is $\zeta_{H_K}$ or the pertinent Hecke factors, not automatically
$\zeta_K$ alone.

\paragraph{Locations.}
The cohomological identification is
\cite[Exp.~IX, Theorem~3.3]{SGA4}; its $K_0$ and $\gamma$ language is in
\cite[Exp.~V 3.10; Exp.~VI 3.2--3.3]{SGA6}.
Cote~2, scan~5 (sheet~4), discusses rank augmentation and Picard comparison
in dimension at most one \cite{Archive2}.  No class-group complexity
theorem is attributed to that manuscript.

\subsection{CX-G04: integral orbit closures and the role of Hilbert and Quot schemes}
\label{finding:CX-G04}

\textbf{Finding.} Let $G=\operatorname{GL}_{m^2,\mathbb Z}$ act on the
coefficient space $V_{\mathbb Z}$ of degree-$m$ forms in $m^2$ variables.
If $I_{\mathbb Q}$ is the ideal of the characteristic-zero determinant
orbit closure, then
\[
 I_{\mathbb Z}=I_{\mathbb Q}\cap\mathbb Z[V],\qquad
 X_{\mathbb Z}=\operatorname{Spec}(\mathbb Z[V]/I_{\mathbb Z})
\]
defines a flat $G$-scheme over $\mathbb Z$.  Its projectivization is a flat
projective family and defines a point of the appropriate Hilbert scheme.
This supplies integral objects for CX-E06--E08 and CX-K05--K08, without
asserting special-fiber equality with every characteristic-$p$ orbit closure.

\begin{proof}
The quotient embeds in $\mathbb Q[V]/I_{\mathbb Q}$ and is torsion free
over $\mathbb Z$.  A torsion-free module over the PID $\mathbb Z$ is a
filtered union of finitely generated torsion-free submodules, which are
free.  Filtered colimits preserve flatness, proving flatness.

The determinant orbit closure is a cone in characteristic zero: over an
algebraic closure, any nonzero scalar multiple of the determinant is
obtained by scalar change of variables.  Its ideal is therefore homogeneous,
as is its contraction.  Every graded quotient piece is finitely generated
and torsion free over $\mathbb Z$, hence free.  On a standard affine chart
of Proj, localization preserves flatness and the degree-zero component is
a direct summand over $\mathbb Z$.  Thus Proj is flat.  It is finitely
presented by noetherianity.  The graded pieces have the same dimensions in
every fiber, so the fiber Hilbert polynomials coincide.  For their polynomial
$P$, \cite[\S3, Theorem~3.2]{FGA221} supplies the classifying morphism
\[
 \operatorname{Spec}\mathbb Z\longrightarrow
 \operatorname{Hilb}^{P}(\mathbb P(V)_{\mathbb Z}/\mathbb Z).
\]

The integral change-of-variables formula defines the original coaction.
For $f\in I_{\mathbb Z}$ its image in
$\mathbb Z[G]\otimes(\mathbb Z[V]/I_{\mathbb Z})$ vanishes after tensoring
with $\mathbb Q$.  This target is flat and torsion free over $\mathbb Z$,
so the image already vanishes.  Hence the ideal is $G$-stable, proving the
asserted integral action.  This argument does not claim that formation of
a scheme-theoretic image commutes with all special-fiber base changes.
\end{proof}

Neither finite presentation nor representability gives an efficient list
of equations or a succinct computation of $P$.  No step bounds the degrees
or bit sizes of ideal generators as a function of the succinct circuit
parameter.  The independent equation-access gate of CX-K06 therefore remains.
The effective statement is the search for a finite morphism in CX-K08;
placing its data in Hilbert or Quot schemes does not compute that search.

\paragraph{Locations.}
Linear spaces, actions, and comodules are
\cite[Exp.~I 4.6.3.1, 4.7.1--4.7.2]{SGA3}; finite morphisms and base change
are \cite[6.1.4--6.1.5]{EGA2}.  The parameter spaces and their projective
embeddings are \cite[\S3, Theorems~3.1--3.2, pp.~259--260;
Proposition~3.8, p.~265]{FGA221}.
The archival companions are cote~28, scan~3 (sheet~2), and cote~98,
scan~44 (sheet~43), headed ``Familles limit\'ees de faisceaux coh\'erents''
\cite{Archive28,Archive98}.  The latter sheet is not used in place of the
published representability theorem.

\subsection{CX-G05: exterior powers, Schur operations, and loss of information under rank}
\label{finding:CX-G05}

\textbf{Finding.} The representation objects and locally free sheaves above
have exact-category Grothendieck rings, tensor multiplication, and exterior
operations $\lambda^i$.  With rank augmentation $\varepsilon$, the usual
operations on rank-zero classes are
\[
 \gamma_t(x)=\lambda_{t/(1-t)}(x),\qquad
 F^j_\gamma=\left\langle
  \prod_a\gamma^{i_a}(x_a):\ \varepsilon(x_a)=0,\ \sum_a i_a\geq j
 \right\rangle .
\]
These are the existing structures of
\cite[Exp.~V 2.4, 3.10(3.10.1); Exp.~VI 3.2--3.3]{SGA6}.
For the Boolean scheme of CX-E07, with $N=\#F$,
\[
 [p_*\mathcal O_{X_F}]=N[\mathcal O],\qquad
 \lambda_t(N[\mathcal O])=(1+t)^N.
\]
Its reduced rank class is zero.  This rank data does not determine the
equivariant coordinate-ring classes required for the obstruction in CX-E08.

\begin{proof}
CX-E07 identifies the pushforward with a direct sum of $N$ copies of
$\mathcal O$.  Its $i$-th exterior power has rank $\binom Ni$, giving the
displayed series.  Subtracting $N[\mathcal O]$ gives zero, so all positive
$\gamma$-operations on that reduced class vanish, functorially after base
change.

For the loss of representation data, the trivial and sign representations
of $C_2$ over $\mathbb Q$ both have rank one.  Their traces at the nonidentity
element are $1$ and $-1$.  Additivity of trace in exact sequences shows that
their $K_0$ classes differ.  Thus rank is not injective.  A count transported
only by the rank-class dictionary cannot specify the equivariant class,
let alone its graded Schur multiplicities.  This refutes that particular
reconstruction from rank; it is not a complexity separation or a claim
that no extra reduction could ever provide such data.
\end{proof}

Over a $\mathbb Q$-scheme, Schur functors on locally free $\mathcal E$ are
images of rational Young idempotents on $\mathcal E^{\otimes d}$.
Idempotents split, giving locally free images.  One can likewise work after
inverting the required factorial denominators.  This is not an assertion
of integral semisimplicity over $\mathbb Z$.
Their classes are universal Schur operations; integrally these classes can
be written using Jacobi--Trudi,
\[
 s_\lambda(\mathcal E)
 =\det(h_{\lambda_i-i+j}(\mathcal E)),\qquad
 \sum_{i\geq0}h_i(\mathcal E)t^i=\lambda_{-t}([\mathcal E])^{-1},
\]
with $h_0=1$ and $h_i=0$ for $i<0$
\cite[Chapter~I \S3, (3.4), p.~41]{Macdonald1995}.
The representation interpretation is the equivariant-module language of
\cite[Exp.~I 4.7.1--4.7.2]{SGA3}.  Cote~2, scan~5 (sheet~4), is the rank
augmentation companion, and cote~98, scan~12 (sheet~11), is the finite
$G$-set companion \cite{Archive2,Archive98}.  No plethysm or Kronecker
problem from the parallel prompts is being asserted solved.

\subsection{CX-G06: the exact boundary between pointwise weights and global GRH}
\label{finding:CX-G06}

\textbf{Finding.} The finite-monodromy arithmetic sheaves in CX-G01 have
root-of-unity Frobenius eigenvalues at good primes, hence pointwise weight
zero.  Tensor, exterior, and characteristic-zero direct-summand operations
preserve that property.  This fact does not supply an additional proof of
the global zero-location input used in Part~I.

\begin{proof}
A finite-order matrix is annihilated by $T^a-1$ for some $a$; all its
eigenvalues are roots of unity.  Tensor eigenvalues and exterior-power
eigenvalues are products of these, and eigenvalues of a stable direct
summand occur among those of the original representation.  All complex
absolute values are one.  This proves precisely the pointwise assertion.

The number-field GRH assertion concerns analytic continuation and zeros of
an infinite Euler product in $s$.  The preceding finite-order linear-algebra
argument supplies no such continuation or zero estimate.  In our proved
implications the required global assertion remains $A_D$ or $M$, exactly
as recorded in CX-G01.  This identifies the scope of the argument, not
logical independence of mathematical theories.

For the Boolean scheme, CX-E07 gives the stronger explicit description.
After base change to $\overline{\mathbb F}_q$, its étale cohomology is $N$
copies of $\mathbb Q_\ell$ in degree zero and zero in higher degrees.
Frobenius is the identity.  Hence the fiber zeta function is $(1-t)^{-N}$
and its Artin motive is a sum of $N$ unit motives.  The counting algorithm
computes this multiplicity from a succinct presentation; it constructs
no nontrivial motive and changes no purity theorem.
\end{proof}

The existing finite-field results are
\cite[\S1.2; Theorem~3.3.1, p.~204; Corollary~3.3.9, p.~207]{WeilII},
with the Euler-factor and direct-image formalism of
\cite[Exp.~XV \S3, no.~1 Proposition~1, and no.~2]{SGA5}.
The inspected archival comparison is cote~34, scan~64 (sheet~63),
crossed-out Remark~1.17, which relates a proposed openness assertion to
Riemann--Weil in higher dimension; scan~74 (sheet~73), equation~(2.5.5),
uses the curve case for an invariant-cohomology vanishing
\cite{Archive34}.  The crossed-out status of the first passage is part
of the locator.  These are finite-field cohomological companions, not
number-field zero-location statements.  Our finding changes neither
their mathematical scope nor the published Weil~II theorem.

\section{Where every finding lives in this dictionary}
\label{sec:crosswalk}

\begin{longtable}{@{}p{0.17\textwidth}p{0.42\textwidth}p{0.34\textwidth}@{}}
\toprule
Findings & Restatement and proof-carrying dictionary & Exact locations and archival companions\\
\midrule
\endhead
CX-R01--R02 & The finite étale permutation object, rank-one decomposition, and commutator obstruction of CX-G01. & SGA1 V \S4 Thm.~1, \S7; SGA3 I4.7.2; cote98/12 (sheet11).\\
CX-K01--K04 & Nonempty generic fiber, bounded closed point, degree-one specializations; abelian promise and cubic obstruction. CX-K08 and CX-G01 carry the proofs. & SGA1 V \S7; SGA5 XV \S3 no.1; cote98/11--12 (sheets10--11).\\
CX-K05--K07 & Search for a finite map from a graded affine scheme to affine space; a small separating system and a dimension-size system retain different output requirements. CX-K08 and CX-G04 preserve these distinctions. & EGAII6.1.4--6.1.5; FGA221 Thms3.1--3.2; cote98/44 (sheet43).\\
CX-K08 & Integral spreading, finite étale cover, Frobenius trace, and open locus of finite linear projections. & SGA1 V \S7; EGAII6.1.4; FGA221 \S3; cote28/3 and cote98/11--12,44.\\
CX-N01--N04 & Explicit residue-ring and $\mathbb G_m$ computations; Kummer torsor classes and generators, retaining the factorization input. CX-G02. & SGA3 VIII1.1--1.2; SGA4 IX3.2; cote28/3, cote98/12.\\
CX-N05--N06 & Ray-character local systems on arithmetic curves and their direct images, together with finite-field group computations. CX-G01--G02. & SGA1 V \S7; SGA3 I4.7.2; SGA4 IX3.2; cote98/11--12.\\
CX-N07--N08, CX-N12 & Picard generators, $S$-units, principal-divisor relations, and the same certification index. CX-G03. & SGA4 IX3.3; SGA6 V3.10, VI3.2; cote2/5 (sheet4).\\
CX-N09--N11, CX-N13--N15 & Connected-component/idempotent computation for finite étale algebras, Kummer power classes, and Frobenius in the stated arithmetic cover. CX-G01--G02; the prime-field automorphism and general composition proofs construct open-and-closed decompositions with their original bounds. & SGA1 V \S7; SGA4 IX3.2; SGA3 VIII1.1; cote98/12, cote28/3.\\
CX-E01--E05 & Emptiness and degree of the succinct Boolean finite étale scheme, parsimonious gate isomorphisms, and the original Boolean circuit model. The table and proof in CX-E07 give each restatement. & SGA1 V \S7; SGA6 VI3.2--3.3; SGA5 XV \S3; cote98/11--12, cote2/5.\\
CX-E06, CX-E08 & Symbolic polynomials and graded equivariant coordinate rings of the integral orbit models; multiplicity quotient condition. CX-G04--G05. & SGA3 I4.7.2; FGA221 Thm3.2; SGA6 V2.4; cote28/3, cote98/44.\\
CX-E07 & Degree equals rank and the multiplicity of the unit Artin motive; reduced $\gamma$-class is zero. CX-G05--G06. & SGA6 V3.10, VI3.2; SGA5 XV \S3; cote2/5, cote98/11.\\
CX-G01--G06 & The displayed geometric proofs are already in the requested language; no further translation changes their premises. & The individual exact locators and scan numbers are supplied with each finding.\\
\bottomrule
\end{longtable}

Here ``cote98/12'' means PDF scan~12 of cote~98, including the archival
cover as scan~1; the sheet itself is numbered~11.  The same convention is
used throughout.  The archive references are to the inspected page content,
not a claim of dating or priority beyond the archive's catalogue.


\clearpage\part{Unproved extensions and verification record}

\section{Unproved extensions: exact status and current gates}
\label{sec:unproved}

This part contains questions and proposed extensions, not additional
theorems. The transfers in Parts~I--II have proofs and recorded checks.
An analytic hypothesis absent from the stipulated theorem is not thereby
proved necessary, and a failed reduction is not a logical independence
result. The following register prevents those distinctions from being
lost when the conclusions are reused.

\paragraph{CX-U01: an unrestricted Dirichlet-only HN transfer.}
The checked general route reaches RH for the Dedekind zeta function of
the bounded solution field (CX-K01). Applying prime-ideal counting to
the field itself avoids the unnecessary Galois-closure gate.
The bounded abelian-witness promise is proved in CX-K03.
The cubic in CX-K04 refutes replacement of every witness by an abelian
one; CX-R02 separately refutes a virtual rank-one representation
decomposition. Neither refutation excludes a different complexity
proof. A new argument relating rational Dirichlet zeros to the needed
prime supply, or a different AM protocol, would be required to go
further. The unconditional prime-ideal estimate inspected in
Lagarias--Odlyzko does not keep the primes at polynomial bit length
with the available solution-field bounds.

\paragraph{CX-U02: an unrestricted Dirichlet-only factoring transfer.}
The main hypothesis classification is proved: CX-N15 gives a complete
general factoring algorithm under $M$, using precisely the sufficient
family $\zeta_{\mathbb Q}$ and
$\zeta_{\mathbb Q(\zeta_{2\nu},p^{1/\nu})}$ for its explicit construction
degrees $\nu\le d^2$. It constructs the maximal coefficient Teichmüller
unit by a Frobenius eigenspace and equivariant interpolation, without
the historical radical oracle or the unretrieved internal EVD94 reduction.
Those retrieval limitations remain historical source facts, not gaps in
CX-N15. CX-N13 and CX-N14 give smaller-hypothesis cases.

The stronger extension asking for a general proof from $A_D$ alone
remains unproved. A reduction to polynomials split over the prime field
does not itself remove the root-of-unity extensions used to split
automorphisms. After CX-N14 the uncovered branch can be restricted to
$p>d$, $r\nmid p(p-1)$ and
$r^2\mid p^{\operatorname{ord}_r(p)}-1$.
A coefficient Teichmüller unit of maximal $r$-power order is a proved
sufficient extra input there. CX-N15 supplies it under its explicit
Dedekind-zeta family; no reduction of that construction to rational
Dirichlet characters, and no impossibility of such a reduction, is proved.
One further natural shortcut was checked directly: BGL Theorem~3
requires roots for every prime divisor of $\Phi_k(p)$ and retains the
largest such prime in its running time, even for one desired nonresidue.
Its alternative Theorem~(5.1) requires a degree-$r$ irreducible polynomial
when $r\nmid k$, exactly the missing input at $k=\operatorname{ord}_r(p)<r$;
see \cite[\S5, p.~20 and algorithm pp.~20--24]{CXNBGL01}.
Thus that source-based substitution does not close the smaller gate.

\paragraph{CX-U03: general polynomial-time finite-field factoring.}
The Evdokimov bound in CX-N09 and the independent composition CX-N15
are quasipolynomial in the degree.
The special-field and structural-parameter results in CX-N10 preserve
their parameters. Deleting an analytic hypothesis changes none of
these exponents. No general deterministic polynomial-time factoring
claim is made. The source-specific transfers answer the requested
hypothesis question without asserting this stronger algorithm.

\paragraph{CX-U04: unrestricted class-group search complexity.}
CX-N07 proves the factor-base transfer, and CX-N08 proves what the
index certificate establishes. The cost of finding enough independent
relations remains an algorithm-specific question. CX-N12 supplies a
concrete rigorous imaginary-quadratic computation theorem with its
randomness and subexponential bound retained. No general smoothness
heuristic is promoted to a theorem by substituting the GRH premise.

\paragraph{CX-U05: equations, symbolic circuits, and GCT separation.}
For an optimal-cardinality succinct h.s.o.p., the historical EXPH
proof still needs the equation-access assumption of GCT~V
Theorem~5.5. A polynomial-size separating system has a different,
already unconditional PSPACE construction. The ETH counter and the
Boolean circuit lower bound supply neither symbolic permanent
circuits nor graded equivariant coordinate-ring multiplicities.
The proved rank dictionary loses information, as CX-G05 demonstrates.
The remaining gates are therefore the specified equation algorithm,
arithmetic conversion, or obstruction construction. No conclusion
about P versus NP, VP versus VNP, or general determinant border
complexity is claimed beyond the explicit statements in Part~I.

\paragraph{CX-U06: stronger quantitative ETH transfers.}
The source theorem fixes clause width. The paper does not interchange
that quantifier with a width growing with input length. A graph's
vertex parameter is not its treewidth or solution-size parameter.
Likewise the Circuit-SAT encoding introduces one variable per gate,
and the nearly linear Boolean lower bound is in
$\mathrm E^{\mathrm{NP}}$. The exact reduction size, solver prefactor,
and computational model must be supplied before any stronger
fine-grained or parameterized conclusion follows.

\section{Source and verification conventions}
\label{sec:source-conventions}

The accompanying \texttt{QUESTIONS.tsv} records each treated question
and its proved finding, refutation, or remaining gate.
\texttt{CHECK\_EVENTS.tsv} records actual checks using the categories
specified in the brief. A targeted proof check means that the stated
proof or the relevant hypothesis-bearing part was read; it does not
mean a formal certificate. A cited theorem whose statement is used
without an independent reproof is identified as such.

The research run is \texttt{CX\_GRH-20261121}, begun at
23:00:34 UTC on 20 November 2026, or 00:00:34 on 21 November in
Europe/Berlin. The coordinator and three researchers used
GPT-6 Astra Pro. Their effort setting was not exposed by the runtime
and is not guessed. Dates in the consolidated registers use the
Europe/Berlin calendar; original UTC dates in component working notes
are explicitly identified. No new finding was checked by a Lean
kernel in this run.

The prior run's \texttt{BRIDGES\_QUESTIONS.tsv} rows Q07 and Q08
were read as earlier work, then rechecked: Q07's broad Dedekind gate
was sharpened to the bounded solution field, and Q08's Dirichlet
scope was verified against Bach's original Theorems~2--4 and
Schoof's complete quadratic-character proof. Their earlier check
labels are not inherited as certificates for this paper.

Archive scan numbers count the cover as scan~1. The paper also gives
the manuscript's sheet number, when present. Mathematical content
was inspected on the cited scans. These historical references locate
constructions and comparisons; none is cited as a manuscript proof
of a modern complexity statement. In particular the crossed-out
passage in cote~34 retains that qualification.

The attached sources remain the primary record for the stipulated
2026 theorems. Literature links in the bibliography identify the
consulted editions or author versions. Differences between printed
journal pagination and author-PDF pagination are stated at the
point of citation. Retrieval failures and source limitations are
recorded without inventing statement numbers or attributing a
full-proof audit to an abstract.


\clearpage\begin{thebibliography}{99}

\bibitem{Alpoge2026} L.~Alp\"oge, \emph{Remark}, 20 November 2026, supplied four-page note,
\url{https://alpo.ge/remark.pdf}; supplied \texttt{ALPOGE\_Main.lean} and
\texttt{ALPOGE\_Statements.lean}. Exact interfaces and step locator in Section~\ref{sec:premises}.


\bibitem{Masullo2026} A.~M.~Masullo, \emph{The Music of the Primes}, supplied 940-page draft,
20 November 2026. Chapter~1 \S1.2; Chapter~3 Definitions~3.2.1.2 and~3.2.1.4,
Proposition~3.2.1.3; Theorems~12.5.1.1 and~13.3.1.1. The empty heading
13.3.1.2 is not used as a theorem statement.


\bibitem{Washington1997} L.~C.~Washington, \emph{Introduction to Cyclotomic Fields},
2nd ed., Graduate Texts in Mathematics 83, Springer, 1997, Theorem~4.3, p.~33.
\url{https://doi.org/10.1007/978-1-4612-1934-7}.


\bibitem{MilneCFT} J.~S.~Milne, \emph{Class Field Theory}, v4.03, 6 August 2020,
Introduction Theorems~0.4--0.5 (Hilbert and ray class fields).
\url{https://www.jmilne.org/math/CourseNotes/CFT.pdf}.


\bibitem{SGA1} A.~Grothendieck et al., \emph{Rev\^etements \'etales et groupe fondamental},
SGA~1, LNM~224, Springer, 1971; Exp.~V \S4, Theorem~1, and \S7.
\url{https://grothendiecksga.com/read/sga1/fr/V-4.html};
\url{https://grothendiecksga.com/read/sga1/fr/V-7.html} (consulted text transcription).


\bibitem{SGA3} M.~Demazure and A.~Grothendieck (eds.), \emph{Sch\'emas en groupes},
SGA~3, LNM~151--153, Springer, 1970; Exp.~I 4.6.3.1, 4.7.1--4.7.2;
Exp.~VIII 1.1--1.2. Revised text consulted at
\url{https://wstein.org/sga/SGA3/Exp1-20jan2010.pdf} and
\url{https://wstein.org/sga/SGA3/Exp8-8nov09.pdf}.


\bibitem{SGA4} M.~Artin, A.~Grothendieck, J.-L.~Verdier (eds.),
\emph{Th\'eorie des topos et cohomologie \'etale des sch\'emas}, SGA~4,
LNM~269, 270, 305, Springer, 1972--1973; Exp.~IX 3.2--3.3.
\url{https://grothendiecksga.com/read/sga4/fr/IX-3.html} (consulted transcription).


\bibitem{SGA5} A.~Grothendieck, with collaborators, \emph{Cohomologie $\ell$-adique
et fonctions $L$}, SGA~5, ed.~L.~Illusie, LNM~589, Springer, 1977;
Exp.~XV \S3, no.~1, Proposition~1, and no.~2, theorem and corollary to Lemma~1.
\url{https://grothendiecksga.com/read/sga5/fr/XV-3.html} (consulted transcription).


\bibitem{SGA6} P.~Berthelot, A.~Grothendieck, L.~Illusie (eds.),
\emph{Th\'eorie des intersections et th\'eor\`eme de Riemann--Roch}, SGA~6,
LNM~225, Springer, 1971; Exp.~V 2.4, 3.10 (3.10.1);
Exp.~VI Theorems~3.2--3.3.
\url{https://grothendiecksga.com/read/sga6/fr/VI-3.html} and
\url{https://grothendiecksga.com/read/sga6/fr/V-3.html} (consulted transcriptions).


\bibitem{EGA2} A.~Grothendieck and J.~Dieudonn\'e, \emph{\'El\'ements de g\'eom\'etrie
alg\'ebrique II}, Publ. Math. IH\'ES 8 (1961), 5--222; 6.1.4--6.1.5.
\url{https://www.numdam.org/item/PMIHES_1961__8__5_0/}.


\bibitem{FGA221} A.~Grothendieck, \emph{Techniques de construction et th\'eor\`emes
d'existence en g\'eom\'etrie alg\'ebrique IV: les sch\'emas de Hilbert},
S\'eminaire Bourbaki, Exp.~221, 1960--1961, pp.~249--276;
\S3 Theorems~3.1--3.2, pp.~259--260, Proposition~3.8, p.~265.
\url{https://www.numdam.org/article/SB_1960-1961__6__249_0.pdf}.


\bibitem{WeilII} P.~Deligne, \emph{La conjecture de Weil II}, Publ. Math. IH\'ES 52
(1980), 137--252; \S1.2, Theorem~3.3.1, p.~204, Corollary~3.3.9, p.~207.
\url{https://publications.ias.edu/sites/default/files/Number40.pdf}.


\bibitem{Macdonald1995} I.~G.~Macdonald, \emph{Symmetric Functions and Hall Polynomials},
2nd ed., Oxford University Press, 1995; Chapter~I \S3, equation~(3.4), p.~41.
\url{https://math.berkeley.edu/~corteel/MATH249/macdonald.pdf}.


\bibitem{Archive2} Archives Grothendieck, Universit\'e de Montpellier, cote~2,
\emph{Riemann--Roch: notes antiques}, scan~5, numbered sheet~4.
\url{https://grothendieck.umontpellier.fr/2.pdf}.


\bibitem{Archive28} Archives Grothendieck, cote~28, \emph{Groupes alg\'ebriques
(SGA~3): notes manuscrites}, circa 1970, scan~3, numbered sheet~2.
\url{https://grothendieck.umontpellier.fr/28.pdf}.


\bibitem{Archive34} Archives Grothendieck, cote~34,
\emph{SGA~7: notes manuscrites}, catalogue dating circa 1967--1973;
scan~64 (sheet~63), crossed-out Remark~1.17; scan~74 (sheet~73),
equation~(2.5.5) and its Riemann--Weil explanation.
\url{https://grothendieck.umontpellier.fr/34.pdf}.


\bibitem{Archive98} Archives Grothendieck, cote~98, scans~11--12 (sheets~10--11),
norm-morphism axioms and finite $G$-sets; scan~44 (sheet~43),
\emph{Familles limit\'ees de faisceaux coh\'erents}.
\url{https://grothendieck.umontpellier.fr/98.pdf}.


\bibitem{K96}
Pascal Koiran, \emph{Hilbert's Nullstellensatz Is in the Polynomial
Hierarchy}, Journal of Complexity \textbf{12} (1996), 273--286.
doi:10.1006/jcom.1996.0019.
Journal copy consulted:
\url{https://cs.uwaterloo.ca/~r5olivei/courses/2025-spring-cs860/K96-HN-PH.pdf}.



\bibitem{K96TR}
Pascal Koiran, \emph{Hilbert's Nullstellensatz Is in the Polynomial
Hierarchy}, DIMACS Technical Report 96-27, July 1996.
Author's extended version:
\url{https://perso.ens-lyon.fr/pascal.koiran/Publis/DIMACS.96-27.ps}.
Archive record:
\url{https://archive.dimacs.rutgers.edu/TechnicalReports/abstracts/1996/96-27.html}.



\bibitem{K97Err}
Pascal Koiran, erratum to \emph{Hilbert's Nullstellensatz Is in the
Polynomial Hierarchy}, 8 August 1997, two pages.
\url{https://perso.ens-lyon.fr/pascal.koiran/Publis/hn.erratum.ps}.



\bibitem{LO77}
J.~C. Lagarias and A.~M. Odlyzko, \emph{Effective Versions of the
Chebotarev Density Theorem}, in A. Fr\"ohlich (ed.), \emph{Algebraic
Number Fields}, Academic Press, 1977, 409--464.
Consulted TeXromancers transcription, Theorem~1.1 on PDF p.~4:
\url{https://aareyanmanzoor.github.io/assets/articles/lagarias-odlyzko.pdf}.
The transcription is not presented as an original scan.



\bibitem{GCT5}
Ketan D. Mulmuley, \emph{Geometric Complexity Theory V: Efficient
algorithms for Noether Normalization}, Journal of the American
Mathematical Society \textbf{30} (2017), 225--309.
doi:10.1090/jams/864.
Attached \texttt{main.tex}; numbered locators checked against
arXiv:1209.5993v5, \url{https://arxiv.org/pdf/1209.5993}.



\bibitem{Burgisser04}
Peter B\"urgisser, \emph{The complexity of factors of multivariate
polynomials}, Foundations of Computational Mathematics \textbf{4}
(2004), 369--396, Theorem~5.7, as invoked in GCT~V Theorem~2.8.



\bibitem{FS18}
Michael A. Forbes and Amir Shpilka, \emph{A PSPACE Construction of a
Hitting Set for the Closure of Small Algebraic Circuits}, Proceedings
of STOC 2018, 1180--1192. doi:10.1145/3188745.3188792.
Consulted ECCC TR17-163, Theorem~1.1 p.~3 and Theorem~7.1 p.~21:
\url{https://eccc.weizmann.ac.il/report/2017/163/download/}.



\bibitem{GSS19}
Zeyu Guo, Nitin Saxena, and Amit Sinhababu,
\emph{Algebraic Dependencies and PSPACE Algorithms in Approximative
Complexity over Any Field}, Theory of Computing \textbf{15}(16)
(2019), 1--30. doi:10.4086/toc.2019.v015a016.
Theorems~1.5 and~1.6, p.~6.
\url{https://www.cse.iitk.ac.in/users/nitin/papers/toc18-ADinPH.pdf}.

% Optional: only if the parent uses the least-prime result elsewhere.


\bibitem{CXNBachSorenson96}
Eric Bach and Jonathan Sorenson, \emph{Explicit bounds for primes in
residue classes}, Mathematics of Computation \textbf{65}(216)
(1996), 1717--1735. doi:10.1090/S0025-5718-96-00763-6.
Theorem~5.1, p.~1729; analytic convention in \S2, p.~1720.
Author-uploaded original text:
\url{https://www.researchgate.net/publication/220576951_Explicit_bounds_for_primes_in_residue_classes}.



\bibitem{BIMS17}
Vishwas Bhargava, G\'abor Ivanyos, Rajat Mittal, and Nitin Saxena,
\emph{Irreducibility and deterministic $r$-th root finding over finite
fields}, Proceedings of ISSAC 2017, eight-page author version.
doi:10.1145/3087604.3087620.
Corollary~1.3 and field-input convention, PDF p.~2.
\url{https://www.cse.iitk.ac.in/users/nitin/papers/irred-and-rth-root.pdf}.
ECCC TR17-016: \url{https://eccc.weizmann.ac.il/report/2017/016/}.




\bibitem{CXBachShallit89}
Eric Bach and Jeffrey Shallit, \emph{Factoring with cyclotomic polynomials},
Mathematics of Computation \textbf{52}(185) (1989), 201--219.
doi:10.1090/S0025-5718-1989-0947467-1.
\S8, Lemma~13 pp.~213--214; Lemma~14 and Theorem~3 p.~214.
Full version of the FOCS~1985 paper, pp.~443--450.
\url{https://www.ams.org/journals/mcom/1989-52-185/S0025-5718-1989-0947467-1/S0025-5718-1989-0947467-1.pdf}.



\bibitem{CXNBach90}
E. Bach, \emph{Explicit bounds for primality testing and related problems},
Math. Comp. \textbf{55} (1990), no.~191, 355--380.
\url{https://doi.org/10.1090/S0025-5718-1990-1023756-8}.
Original journal locators: \S6, pp.~369--373, Theorems~2--3 p.~372;
Miller consequence p.~373; Theorem~4 p.~375; class-group corollary p.~376.



\bibitem{CXNSchoof08}
R. Schoof, \emph{Four primality testing algorithms}, in
\emph{Algorithmic Number Theory: Lattices, Number Fields, Curves and Cryptography},
MSRI Publications \textbf{44} (2008), 101--126.
Theorem~2.2, pp.~105--106.
\url{https://pub.math.leidenuniv.nl/~stevenhagenp/ANTproc/05rene.pdf}.



\bibitem{CXNMiller76}
G. L. Miller, \emph{Riemann's hypothesis and tests for primality},
J. Comput. System Sci. \textbf{13} (1976), 300--317.
Theorem~2 p.~301; Appendix pp.~315--316.
\url{https://www.cs.cmu.edu/~glmiller/Publications/Papers/Mi76.pdf}.



\bibitem{CXNShoup92}
V. Shoup, \emph{Searching for primitive roots in finite fields},
Math. Comp. \textbf{58} (1992), no.~197, 369--380.
\url{https://doi.org/10.1090/S0025-5718-1992-1106981-9}.
Locators in this section refer to the 16-page author preprint:
Theorems~2--3 pp.~3--4, Proposition~9 p.~10, Proposition~13 pp.~11--12,
proofs \S\S5--6 pp.~12--13.
\url{https://www.shoup.net/papers/primroots.pdf}.



\bibitem{CXNBuchmannShoup96}
J. Buchmann and V. Shoup,
\emph{Constructing nonresidues in finite fields and the Extended Riemann Hypothesis},
Math. Comp. \textbf{65} (1996), no.~215, 1311--1326.
Main statement pp.~1311--1312; ERH definition and Fact~2.1 p.~1314;
Fact~4.3 p.~1319; algorithm \S5.
\url{https://shoup.net/papers/nonresidues.pdf}.



\bibitem{CXNAdlemanLenstra86}
L. M. Adleman and H. W. Lenstra, Jr.,
\emph{Finding irreducible polynomials over finite fields},
Proc. 18th ACM Symposium on Theory of Computing (1986), 350--355.
Theorem~1 p.~350; Algorithm~A pp.~350--351; Proposition~3 p.~351;
Lemma~6 p.~352.
\url{https://pub.math.leidenuniv.nl/~lenstrahw/PUBLICATIONS/1986a/art.pdf}.



\bibitem{CXNBDF08}
K. Belabas, F. Diaz y Diaz, and E. Friedman,
\emph{Small generators of the ideal class group},
Math. Comp. \textbf{77} (2008), no.~262, 1185--1197.
\url{https://doi.org/10.1090/S0025-5718-07-02003-0}.
Locators refer to the author version dated 5 December 2005:
introduction pp.~2--3, Theorem~1.1 p.~3, Theorem~2.1 p.~5,
proof pp.~4--6.
\url{https://www.math.u-bordeaux.fr/~kbelabas/research/OnBach.pdf}.



\bibitem{CXNLenstra92}
H. W. Lenstra, Jr., \emph{Algorithms in algebraic number theory},
Bull. Amer. Math. Soc. \textbf{26} (1992), no.~2, 211--244.
Locators refer to the author preprint:
\S5 following Problem~5.6, exact sequence and (5.7), p.~25.
\url{https://pub.math.leidenuniv.nl/~lenstrahw/PUBLICATIONS/1992a/art.pdf}.



\bibitem{CXNBiasseErukulangara}
J.-F. Biasse and M. R. Erukulangara,
\emph{A proof of the conjectured run time of the Hafner--McCurley class group algorithm},
Adv. Math. Commun. \textbf{17} (2023), no.~6, 1388--1405.
\url{https://doi.org/10.3934/amc.2021055}.
Locators refer to the accepted author PDF: Main Theorem and Conjecture~1 p.~4;
Proposition~1 p.~8; Algorithm~4 p.~15; Proposition~6 pp.~15--16;
Corollary~1 p.~16.
\url{https://par.nsf.gov/servlets/purl/10324557}.



\bibitem{CXNJMV09}
D. Jao, S. D. Miller, and R. Venkatesan,
\emph{Expander graphs based on GRH with an application to elliptic curve cryptography},
J. Number Theory \textbf{129} (2009), 1491--1504.
\url{https://doi.org/10.1016/j.jnt.2008.11.006}.
Theorem~1.1 pp.~3--4, Corollary~1.3 p.~4, proofs pp.~9--10
in arXiv:0811.0647v2.
\url{https://arxiv.org/pdf/0811.0647}.



\bibitem{CXNEvdokimov94}
S. A. Evdokimov,
\emph{Factorization of polynomials over finite fields in subexponential time under GRH},
Algorithmic Number Theory, LNCS \textbf{877} (1994), 209--219.
\url{https://doi.org/10.1007/3-540-58691-1_58}.
Runtime and radical-oracle formulation: abstract p.~209.
The complete algebraic proof was not independently retrieved for this audit.



\bibitem{CXNEvdokimov89}
S. A. Evdokimov,
\emph{Factoring a solvable polynomial over a finite field and Generalized Riemann Hypothesis},
Zap. Nauchn. Sem. LOMI \textbf{176} (1989), 104--117;
English translation J. Soviet Math. \textbf{59} (1992), no.~3, 842--849.
\url{https://doi.org/10.1007/BF01104107}.
Russian-original locators: character convention p.~105;
Theorem~I.1 pp.~107--108; Theorem~I.2 pp.~108--109;
Theorem~I.3 pp.~109--110; Theorem~2.1 p.~110.
\url{https://www.mathnet.ru/eng/znsl4535}.



\bibitem{CXNGao01}
S. Gao, \emph{On the deterministic complexity of factoring polynomials},
J. Symbolic Comput. \textbf{31} (2001), 19--36.
\url{https://doi.org/10.1006/jsco.1999.1001}.
Locators refer to the author version dated 17 March 1999:
prime-field reduction pp.~1--2; \S2.4 and Lemma~2.6 pp.~7--9;
Lemma~3.6 p.~16.
\url{https://www.math.clemson.edu/~sgao/papers/FacPoly.pdf}.



\bibitem{CXNBGL01}
E. Bach, J. von zur Gathen, and H. W. Lenstra, Jr.,
\emph{Factoring polynomials over special finite fields},
Finite Fields Appl. \textbf{7} (2001), 5--28.
\url{https://doi.org/10.1006/ffta.2000.0306}.
Theorem~1 and corollary p.~6; Theorems~2--3 p.~7.
\url{https://pub.math.leidenuniv.nl/~lenstrahw/PUBLICATIONS/2001a/art.pdf}.



\bibitem{CXNRonyai92}
L. R\'onyai, \emph{Galois groups and factoring polynomials over finite fields},
SIAM J. Discrete Math. \textbf{5} (1992), no.~3, 345--365.
\url{https://doi.org/10.1137/0405026}.
The parameterized statement was checked in the primary abstract;
the complete original proof was not independently retrieved.



\bibitem{CXNIKRS08}
G. Ivanyos, M. Karpinski, L. R\'onyai, and N. Saxena,
\emph{Trading GRH for algebra: algorithms for factoring polynomials and related structures},
ECCC Report TR08-099 (2008).
Main Theorem p.~3; \S4.1--\S4.2, Lemma~4.2 p.~14,
Lemma~4.5 p.~16 and Proposition~4.6 pp.~16--17
in the report's printed numbering.
\url{https://eccc.weizmann.ac.il/report/2008/099/download/}.



\bibitem{CXNLenstra91}
H. W. Lenstra, Jr., \emph{Finding isomorphisms between finite fields},
Math. Comp. \textbf{56} (1991), no.~193, 329--347.
\url{https://doi.org/10.1090/S0025-5718-1991-1052099-2}.
Theorem~(1.2), p.~330; explicit-data convention pp.~329--330.
\url{https://pub.math.leidenuniv.nl/~lenstrahw/PUBLICATIONS/1991b/art.pdf}.



\bibitem{CXNAKS04}
M. Agrawal, N. Kayal, and N. Saxena, \emph{PRIMES is in P},
Ann. of Math. \textbf{160} (2004), no.~2, 781--793.
\url{https://doi.org/10.4007/annals.2004.160.781}.
Theorem~4.1 p.~785 and Theorem~5.1 p.~789.
\url{https://annals.math.princeton.edu/wp-content/uploads/annals-v160-n2-p12.pdf}.


\bibitem{OpenAIETH}
OpenAI, \emph{The Exponential Time Hypothesis is False},
OpenAI Math Release preprint, 17 October 2026, revised 25 October 2026.
Supplied LaTeX source:
\texttt{LIT\_OpenAI\_ETH\_is\_False/main\_full\_version.tex};
Theorem~1.1 (\texttt{thm:main}), Corollaries~1.3 and~1.4,
\S1.1.1--\S1.1.3, Appendix~C.1; the quantitative circuit proof is
\texttt{quantitative-circuits.tex}.
\url{https://github.com/openai/math/blob/main/preprints/The-Exponential-Time-Hypothesis-is-False-October-25-2026/paper.pdf}.



\bibitem{Williams2013}
R. Williams, \emph{Improving exhaustive search implies superpolynomial
lower bounds}, SIAM J. Comput. \textbf{42} (2013), 1218--1244.
\url{https://doi.org/10.1137/10080703X}.
Theorem numbering used here follows the author's full version dated
4 May 2010: Theorems~1.1 (printed p.2), 2.1 (p.7), and~6.1
(pp.20--21).
\url{https://people.csail.mit.edu/rrw/improved-algs-lbs2.pdf}.



\bibitem{JMV2013}
H. Jahanjou, E. Miles, and E. Viola, \emph{Local reductions},
ECCC Report TR13-099, revision~2, 27 November 2013,
Theorem~1 (printed p.2), Corollaries~2--3 (p.3), \S6.
\url{https://eccc.weizmann.ac.il/report/2013/099/revision/2/download/}.
Conference version: ICALP 2015, 749--760.



\bibitem{Valiant1979}
L. G. Valiant, \emph{Completeness classes in algebra},
Proceedings of STOC 1979, 249--261; Theorem~2, p.254.
\url{https://doi.org/10.1145/800135.804419}.
Inspected scan:
\url{https://sites.math.washington.edu/~billey/colombia/references/valiant.completeness.1979.pdf}.



\bibitem{Burgisser2024}
P. B\"urgisser, \emph{Completeness classes in algebraic complexity theory},
arXiv:2406.06217v1, 10 June 2024.
Definitions~2.9 (p.6), 2.25 (p.12);
Theorems~2.20 (p.10), 2.29 (p.13), and~7.3 (p.37).
\url{https://arxiv.org/pdf/2406.06217}.


\end{thebibliography}

\end{document}